\chapter{Rewriting}\label{ch:rewriting}
\sloppy

The package \module{archytaszx.rewrite} turns a rule of the ZXW calculus into an operation on a
\py!Diagram! (\cref{ch:diagrams}). It holds five layers. The \emph{rule layer} defines what a rule
is: a pattern, a builder, side conditions, quantifiers, dimension guards, and an exact scalar
(\cref{sec:rewriting-rule}). The \emph{matching layer} locates occurrences of a pattern: a family of
fixed-shape matchers (\cref{sec:rewriting-fixed}), a generic subgraph matcher driven by a pattern
diagram (\cref{sec:rewriting-template}), matchers for the W-node and multiplier rules of qufinite
ZXW (\cref{sec:rewriting-zxw}), and matchers for port-scoped \bboxes{}
(\cref{sec:rewriting-box}). The \emph{engine layer} applies one match without mutating its input,
records provenance, and runs strategies (\cref{sec:rewriting-engine}). The \emph{search layer}
computes normal forms (\cref{sec:rewriting-nf}), saturates an e-graph of diagrams
(\cref{sec:rewriting-egraph}), runs tactics and bidirectional proof search
(\cref{sec:rewriting-tactics}), and rewrites the sheets of a scalable diagram
(\cref{sec:rewriting-sheet}). The \emph{cache layer} memoizes digests, matches, and derived values
across all of these (\cref{sec:rewriting-cache}).

The rules themselves are defined elsewhere: the rule library and lookup in \cref{ch:rules}, the
complete axiom set in \cref{ch:completeness}, and the discard rules and exhaustive search in
\cref{ch:mixed}. The calculus follows Coecke and Duncan~\cite{coecke2011interacting} for ZX,
Poór et al.~\cite{poor2023zxw} and Wang et al.~\cite{wang2021qufinite,wang2023qufinitezxw} for
qudit and qufinite ZXW, and Dixon, Duncan, and Kissinger~\cite{dixon2009graphical,kissinger2014pattern}
for \bboxes{} and their matching.

\begin{table}[ht]
\centering\small
\caption{Modules of the rewriting layer.}\label{tab:rewriting-modules}
\begin{tabularx}{\textwidth}{@{}lX@{}}
\toprule
Module & Content \\
\midrule
\file{rule.py} & \py!Rule!, \py!Pattern!, \py!Match!, \py!BuildResult!, side conditions, dimension constraints and guards, errors. \\
\file{match.py} & Fusion matcher and the fixed-shape matchers: $F^4$, cap, identity, triangle pair, state copy, Hopf, bialgebra, $F$ on a state. \\
\file{match_template.py} & Generic subgraph matching against a pattern diagram, and the rebuild of a right-hand side. \\
\file{match_zxw.py} & Matchers for W fusion, W unit, multiplier rules, and the Z--W bialgebra family. \\
\file{match_box.py} & Matchers for rules that cross a port-scoped \bbox{} boundary. \\
\file{engine.py} & \py!apply!, \py!RewriteStep!, strategies, the termination guard. \\
\file{normal_form.py} & The normal-form driver, loop detection, keys, and comparison views. \\
\file{egraph.py} & The diagram-level e-graph, its scheduler, extraction, and explanations. \\
\file{tactics.py} & The tactic language and bidirectional proof search. \\
\file{sheet.py} & Sheet rewrites of scalable diagrams. \\
\file{cache.py} & Digests, fingerprints, the LRU cache, the match cache, value memos, incremental rematching. \\
\bottomrule
\end{tabularx}
\end{table}

\paragraph{Errors.}
Every error raised in the package derives from \py!RewriteError!, a subclass of Python's
\py!Exception!. Two subclasses partition the failures.
\begin{itemize}
\item \py!RewriteDomainError!: a value or state lies outside the mathematical domain of the
  operation. Raised for a match whose side-condition outcomes do not exactly cover, or do not all
  pass, its rule's conditions; a builder whose introduced scalar disagrees with its rule's; a
  fusion match that disagrees with a fresh resolution; a phase whose entries fall outside the
  shared dimension; a wire or boundary entry naming a consumed port absent from the port mapping;
  and a rewrite that introduces a hard validation issue the input did not carry.
\item \py!RewriteGrammarError!: a request is malformed independent of mathematics. Raised for a
  value object built outside its contract, a match handed to a builder that does not take its
  type, an unevaluable fusion request (equal node ids, an absent node, a wire not in the diagram
  or not incident on both nodes, an out-of-range port index), a malformed \py!BuildResult!, a port
  mapping that collapses a wire, and an unknown rule name at lookup.
\end{itemize}

%======================================================================
\section{Rules}\label{sec:rewriting-rule}

The module \module{archytaszx.rewrite.rule} defines the value objects of a rewrite. None of them
performs graph surgery. A \term{rule} is a frozen record whose \emph{pattern} locates occurrences
in a diagram, whose \emph{builder} constructs the right-hand side at one occurrence, and whose
metadata (side conditions, quantifiers, dimension guards, scalar) states the conditions and the
exact scalar of the equation
\begin{equation}\label{eq:rewriting-rule}
  \sem{D} \;=\; s \cdot \sem{D'} ,
\end{equation}
where $D$ is the matched diagram, $D'$ the rewritten one, and $s$ the scalar the rule introduces.
The engine (\cref{sec:rewriting-engine}) multiplies $s$ into the diagram's global scalar, so the
rewritten diagram denotes exactly what the original denotes.

\subsection{Side conditions}

\begin{definition}[Side condition]\label{def:rewriting-sidecond}
A \term{side condition} is a frozen record \py!SideCondition(name: str, description: str)!. A
pattern declares an ordered tuple of side conditions once. A \term{side-condition outcome} is a
frozen record
\begin{signature}
\py!SideConditionOutcome(name: str, passed: bool, detail: str, deferred: bool = False)!
\end{signature}
holding the result of one named condition against one candidate. The flag \py!deferred! is true
for an outcome that passed on an assumed rather than a verified dimension equality.
\end{definition}

A pattern may declare a condition that always passes; it then records a structural fact for the
certificate. Equality and hashing of both records are field-wise.

\begin{signature}
\py!check_side_condition_coverage(match: Match, side_conditions: tuple[SideCondition, ...], context: str) -> None!
\end{signature}
\begin{algorithm}[Coverage check]\label{alg:rewriting-coverage}
\leavevmode
\begin{enumerate}
\item Collect the outcome names of \py!match.side_condition_outcomes! in order; record every name
  that occurs twice as a duplicate.
\item Let \emph{missing} be the declared names absent from the outcomes and \emph{unexpected} the
  outcome names not declared.
\item If any of the three sets is nonempty, raise \py!RewriteDomainError! with the message
  \code!<context>: side_condition_outcomes do not exactly cover the declared side conditions (duplicate=[..], missing=[..], unexpected=[..])!, each list sorted.
\item Otherwise, if any outcome failed, raise \py!RewriteDomainError! with
  \code!<context>: match has failing side condition(s) [names]!, the names in outcome order.
\end{enumerate}
\end{algorithm}
Both \py!apply! and each builder call this check first, since a builder can be invoked directly.

\subsection{Dimension constraints}

A rewrite between legs of symbolic dimension can rest on an equality that dimension unification
(\cref{ch:algebra}) does not prove outright. Such an equality is recorded, never silently
accepted. \py!Dim.unify! returns one of three outcomes: \emph{success} (possibly with bindings of
free symbols), \emph{deferred} (undecided), or \emph{failure}.

\begin{definition}[Constraint source]\label{def:rewriting-csource}
A \term{constraint source} names the check that produced a dimension constraint. It is a frozen
record \py!ConstraintSource(kind, port_ref=None, node_id=None)! with \py!kind! one of three
\py!ConstraintSourceKind! members:
\begin{itemize}
\item \code!CONNECTING_PAIR! (value \code!"connecting_pair"!): the two ports a consumed wire joins,
  related to each other. Both references are \py!None!.
\item \code!SURVIVING_LEG! (\code!"surviving_leg"!): one surviving port, named by \py!port_ref!,
  related to the running shared dimension. Exactly \py!port_ref! is set.
\item \code!NODE_PHASE! (\code!"node_phase"!): the phase vector of the node \py!node_id!, related to
  the running shared dimension. Exactly \py!node_id! is set.
\end{itemize}
Any other combination raises \py!RewriteGrammarError! at construction. The class methods
\py!connecting_pair()!, \py!surviving_leg(port_ref)!, and \py!node_phase(node_id)! build the three
kinds. The string form is \code!connecting pair!, \code!surviving leg <node>.<direction>[<index>]!,
or \code!phase on node <node>!.
\end{definition}

\begin{definition}[Dimension constraint]\label{def:rewriting-dconstraint}
A \term{dimension constraint} is a frozen record
\begin{signature}
\py!DimensionConstraint(assumed: Dim, equal_to: Dim, source: ConstraintSource, outcome: ConstraintOutcome, bound_here: tuple[tuple[str, Dim], ...] = ())!
\end{signature}
stating one equality \py!assumed! $=$ \py!equal_to! that a rewrite assumed rather than verified as a
syntactic identity. \py!ConstraintOutcome! has two members: \code!DEFERRED! (\code!"deferred"!),
when unification could not decide the equality, and \code!BOUND! (\code!"bound"!), when it
succeeded only by binding a free symbol. A unification that is a bare syntactic identity is never
recorded. The field \py!bound_here! holds exactly the bindings this check produced, sorted by
symbol name; it may bind a symbol to a non-concrete dimension.
\end{definition}

Construction raises \py!RewriteGrammarError! when \py!outcome! is \code!BOUND! and \py!bound_here!
is empty, when \py!outcome! is \code!DEFERRED! and \py!bound_here! is nonempty, or when
\py!bound_here! is not a tuple of \py!(str, Dim)! pairs. The property \py!deferred! is true iff the
outcome is \code!DEFERRED!. The string form is \code!<assumed> == <equal_to> (<outcome>, <source>)!.
The property \py!equation! renders both sides through \tool{SymPy}{https://www.sympy.org}~\cite{meurer2017sympy} and returns \code!a = b!, where the
sides are ordered by the key (side is an integer literal, length, text); a symbolic side therefore
comes first, and of two symbolic sides the shorter.

\begin{signature}
\py!assumed(constraints: Iterable[DimensionConstraint]) -> tuple[DimensionConstraint, ...]!
\end{signature}
returns the constraints that are not identities, one per equation. Two constraints state the same
equation when the unordered pair of their sides, each expanded by \py!sympy.expand!, is equal. A
constraint whose two expanded sides coincide is dropped. The first occurrence of each equation
survives, in input order.

\subsection{Dimension guards}

\begin{definition}[Dimension guard]\label{def:rewriting-guard}
A \term{dimension guard}
\py!DimensionGuard(kind: DimensionGuardKind, value: int | None = None)! is an arithmetic restriction on the dimension $d$ a match resolves to. Its kinds are
\code!PRIME!, \code!COMPOSITE!, \code!AT_LEAST!, \code!EQUALS!, and \code!DIVISIBLE_BY!. The last
three require \py!value!; the first two forbid it. A value must be a Python \py!int! (never a
\py!bool!) and at least $1$. Violations raise \py!RewriteGrammarError!.
\end{definition}

\py!DimensionGuard.evaluate(dim: Dim | None) -> GuardOutcome! returns \code!UNDECIDED! when
\py!dim! is \py!None! or not concrete. Otherwise, with $v$ the integer value of \py!dim!, it returns
\code!PASSED! or \code!FAILED! according to the kind:
\begin{itemize}
\item \code!PRIME!: $v$ is prime (SymPy \py!isprime!);
\item \code!COMPOSITE!: $v \ge 2$ and $v$ is not prime;
\item \code!AT_LEAST!: $v \ge \mathit{value}$;
\item \code!EQUALS!: $v = \mathit{value}$;
\item \code!DIVISIBLE_BY!: $v \bmod \mathit{value} = 0$.
\end{itemize}
Thus $d = 1$ fails both \code!PRIME! and \code!COMPOSITE!. The string forms are
\code!d is prime!, \code!d is composite!, \code!d >= v!, \code!d == v!, and \code!d divisible by v!.

\subsection{Quantifiers}

\py!Quantifiers(leg_counts: tuple[str, ...] = (), dimensions: tuple[str, ...] = ())! names the
port-count variables and the dimension variables over which a rule's equation is stated. The
rewriting layer stores this metadata and does not consult it.

\subsection{Matches, patterns, and builders}

\begin{definition}[Match]\label{def:rewriting-match}
A \term{match} is any object that exposes four read-only properties:
\py!side_condition_outcomes!, the outcomes of every condition the pattern checked;
\py!dimension_constraints!, at most one constraint per constraint source;
\py!all_side_conditions_passed!, true iff every outcome passed (vacuously true for an empty
tuple); and \py!support_node_ids!, every node id the match reads or rewrites, ascending and
without repetition. \py!Match! is a structural protocol: each pattern defines its own frozen
match type and carries its location data there. The engine never reads that data.
\end{definition}

\begin{definition}[Pattern]\label{def:rewriting-pattern}
A \term{pattern} is an instance of the abstract class \py!Pattern! with
\begin{itemize}
\item \py!locality_radius: int!, default $3$: how many wire hops beyond a match's support the
  pattern's decision can read;
\item \py!find_matches(diagram) -> tuple[Match, ...]! (abstract): every occurrence, in a
  deterministic order; it never mutates the diagram and never returns a candidate with a failed
  side condition;
\item \py!order_key(match) -> tuple!: the key by which \py!find_matches! orders its result; the
  default is the tuple of integer support node ids;
\item \py!verify(diagram, match) -> bool!: by default, whether \py!find_matches! rediscovers the
  match;
\item \py!find_matches_anchored(diagram, anchors: frozenset[NodeId])!: by default, the matches
  whose support meets \py!anchors!. Subclasses override it to restrict the search itself.
\end{itemize}
\end{definition}

The invariant \py!find_matches(D) == tuple(sorted(find_matches(D), key=order_key))! holds for
every pattern in the package. The anchored search returns the same matches as the full search
filtered to the anchors' neighborhood, and raises what the full search raises.

\begin{definition}[Build result]\label{def:rewriting-buildresult}
A \term{builder} is a callable \py!RuleBuilder = Callable[[Diagram, Match], BuildResult]!. It
receives the engine's working copy and a match, adds the replacement nodes to that copy in place,
and returns a frozen record
\begin{signature}
\py!BuildResult(diagram: Diagram, new_node_ids: tuple[NodeId, ...], consumed_node_ids: tuple[NodeId, ...], consumed_wires: tuple[Wire, ...], port_mapping: Mapping[PortRef, PortRef], scalar_introduced: Scalar, new_wires: tuple[Wire, ...] = (), phase_substitutions: Mapping[NodeId, Mapping[str, Dim]] | None = None)!
\end{signature}
\begin{itemize}
\item \py!diagram! is the very working object the builder received (checked by identity).
\item \py!new_node_ids! lists every created node, in a deterministic order.
\item \py!consumed_node_ids! and \py!consumed_wires! name what the rewrite removes.
\item \py!port_mapping! sends each port of a consumed node that some surviving wire or boundary
  entry uses to its replacement port.
\item \py!new_wires! joins new nodes to one another; the engine adds them after remapping.
\item \py!phase_substitutions! records, per consumed node, the bindings the builder substituted
  into a phase; \py!None! means none.
\end{itemize}
A builder never removes matched nodes, never edits the wire set, and never edits boundaries; the
engine does all three.
\end{definition}

\subsection{The rule record}

\begin{definition}[Rule]\label{def:rewriting-rulerecord}
A \py!Rule! is the frozen record
\begin{signature}
\py!Rule(name: str, pattern: Pattern, builder: RuleBuilder, side_conditions: tuple[SideCondition, ...], quantifiers: Quantifiers, scalar_introduced: Scalar, scalar_in_dim: Callable[[Dim], Scalar] | None = None, dimension_guards: tuple[DimensionGuard, ...] = (), scalar_in_match: Callable[[Match], Scalar] | None = None, automatic: bool = True)!
\end{signature}
\py!name! is the stable identifier a certificate cites. \py!scalar_introduced! is the exact scalar
(\cref{ch:algebra}) of every application, written in the rule's own dimension symbol when it
depends on the dimension. \py!automatic! states whether the automatic layers (normal form,
saturation, tactics) may use the rule.
\end{definition}

\paragraph{Scalar of an application.}
\py!scalar_for(dim)! returns \py!scalar_in_dim(dim)! when that function is set and \py!dim! is not
\py!None!, and \py!scalar_introduced! otherwise. When \py!scalar_in_match! is set, \py!scalar_for!
raises \py!RewriteGrammarError! (the scalar depends on a matched port count, not on the dimension
alone). \py!scalar_for_match(match)! returns \py!scalar_in_match(match)! when that function is set,
and otherwise \py!scalar_for(match.shared_dim)!, with \py!None! when the match has no
\py!shared_dim! attribute.

\paragraph{Validation.}
Construction raises \py!RewriteGrammarError! when any of these holds:
\py!name! is not a nonempty string; \py!pattern! is not a \py!Pattern!; \py!builder! is not
callable; \py!side_conditions! is not a tuple of \py!SideCondition!; a side-condition name occurs
twice (message lists the repeated names, sorted); \py!quantifiers! is not a \py!Quantifiers!;
\py!scalar_introduced! is not a \py!Scalar! (a \py!float! is rejected); \py!dimension_guards! is not
a tuple of \py!DimensionGuard!; \py!scalar_in_dim! or \py!scalar_in_match! is neither \py!None! nor
callable; \py!automatic! is not a \py!bool!; the builder carries an attribute
\py!side_conditions! whose tuple differs from the rule's; or the pattern carries an attribute
\py!dimension_guards! whose tuple differs from the rule's. The last two checks keep exactly one
source of truth: a builder or pattern that declares its own conditions or guards must agree with
the wrapping rule.

%======================================================================
\section{Fixed-shape matchers}\label{sec:rewriting-fixed}

The module \module{archytaszx.rewrite.match} implements the matchers whose shape is fixed in code.
Each matcher is a module-level function \py!find_..._matches(diagram, *, anchors=None, ...)!
together with a \py!Pattern! subclass that delegates to it, passing \py!anchors! from
\py!find_matches_anchored!. Every matcher reads a diagram and never mutates it. Unless stated
otherwise a matcher does not require a well-formed diagram and never calls the validator.
\Cref{fig:rewriting-fixed-shapes} shows the shapes the matchers locate.


\begin{figure}[htbp]
\centering
\begin{minipage}[b]{0.23\textwidth}\centering
\begin{tikzpicture}[zxd]
  \node[fbox] (f1) at (0,0.6) {$F$}; \node[fbox] (f2) at (0,1.3) {$F$};
  \node[fbox] (f3) at (0,2.0) {$F$}; \node[fbox] (f4) at (0,2.7) {$F$};
  \draw[wire] (0,0) -- (f1) -- (f2) -- (f3) -- (f4) -- (0,3.3);
  \node at (0.7,1.65) {$=$}; \draw[wire] (1.2,0) -- (1.2,3.3);
\end{tikzpicture}\\(a) $F^4$ cancellation
\end{minipage}\hfill
\begin{minipage}[b]{0.23\textwidth}\centering
\begin{tikzpicture}[zxd]
  \node[zdot] (s) at (0,0.6) {}; \node[xdot] (e) at (0,1.8) {};
  \draw[wire] (s) -- (e);
  \node[dimlabel,right] at (0.05,1.2) {$d$};
\end{tikzpicture}\\(b) cap: Z state into X effect
\end{minipage}\hfill
\begin{minipage}[b]{0.23\textwidth}\centering
\begin{tikzpicture}[zxd]
  \node[zdot] (z) at (0,1.2) {};
  \draw[wire] (0,0) -- (z) -- (0,2.4);
  \node at (0.6,1.2) {$=$}; \draw[wire] (1.1,0) -- (1.1,2.4);
\end{tikzpicture}\\(c) identity removal
\end{minipage}\hfill
\begin{minipage}[b]{0.23\textwidth}\centering
\begin{tikzpicture}[zxd]
  \node[tri] (t) at (0,0.8) {}; \node[trii] (ti) at (0,1.7) {};
  \draw[wire] (0,0) -- (t) -- (ti) -- (0,2.5);
  \node at (0.6,1.25) {$=$}; \draw[wire] (1.1,0) -- (1.1,2.5);
\end{tikzpicture}\\(d) $T$ then $T^{-1}$
\end{minipage}

\bigskip
\begin{minipage}[b]{0.23\textwidth}\centering
\begin{tikzpicture}[zxd]
  \node[xdot] (s) at (0,0.4) {}; \node[zs] (z) at (0,1.4) {};
  \draw[wire] (s) -- (z);
  \draw[wire] (z) -- (-0.5,2.3); \draw[wire] (z) -- (0.5,2.3);
  \node at (0,2.1) {$\cdots$};
\end{tikzpicture}\\(e) state copy
\end{minipage}\hfill
\begin{minipage}[b]{0.23\textwidth}\centering
\begin{tikzpicture}[zxd]
  \node[zs] (z) at (0,0.3) {}; \node[xs] (x) at (0,2.6) {};
  \node[fbox] (f1) at (0.6,1.05) {$F$}; \node[fbox] (f2) at (0.6,1.85) {$F$};
  \draw[wire] (z) to[bend left=25] (x);
  \draw[wire] (z) -- (f1) -- (f2) -- (x);
\end{tikzpicture}\\(f) Hopf pattern
\end{minipage}\hfill
\begin{minipage}[b]{0.23\textwidth}\centering
\begin{tikzpicture}[zxd]
  \node[xdot] (x) at (0,0.8) {}; \node[zdot] (z) at (0,1.7) {};
  \draw[wire] (-0.4,0) -- (x); \draw[wire] (0.4,0) -- (x);
  \draw[wire] (x) -- (z);
  \draw[wire] (z) -- (-0.4,2.5); \draw[wire] (z) -- (0.4,2.5);
\end{tikzpicture}\\(g) bialgebra pattern
\end{minipage}\hfill
\begin{minipage}[b]{0.23\textwidth}\centering
\begin{tikzpicture}[zxd]
  \node[zdot] (s) at (0,0.4) {}; \node[fbox] (f) at (0,1.3) {$F$};
  \draw[wire] (s) -- (f) -- (0,2.2);
\end{tikzpicture}\\(h) $F$ on a Z state
\end{minipage}
\caption{Left-hand shapes of the fixed-shape matchers of \file{match.py}, inputs at the bottom.
Small dots are phaseless spiders. Each swapped variant exchanges the two colors.}
\label{fig:rewriting-fixed-shapes}
\end{figure}

\subsection{Shared predicates}\label{sec:rewriting-fixed-shared}

The following predicates recur in this section and in \cref{sec:rewriting-template,sec:rewriting-zxw}.
\begin{itemize}
\item \emph{Registered spider}: a node whose generator type is registered in the generator
  registry and is named \code!Z! or \code!X! (\cref{ch:diagrams}).
\item \emph{Phaseless}: the node has no phase vector, or an all-zero one.
\item \emph{Uniform port dimension}: the set of the dimensions of all ports of the node has exactly
  one element; a node without ports has none.
\item \emph{Claims}: the claim count of a port is the number of wires with that port as an
  endpoint. A port is \emph{claimed exactly once by a wire} when its count is $1$ and no boundary
  list names it; it is \emph{claimed at most once} when its count plus $1$ if a boundary list
  names it is at most $1$.
\item \emph{Node boxes}: the set of node-scoped \bboxes{} whose node scope contains a node.
  The \term{innermost node-scoped box} of a node is the node-scoped \bbox{} containing it with the
  smallest node scope, ties broken by the smaller box id; \py!None! when there is none.
\item \emph{Outside every \bbox{}}: no node of a set lies in a node-scoped \bbox{}, no port of
  such a node lies in a port-scoped \bbox{}, and no wire joins such a node to a node whose node
  boxes differ (a \emph{crossing leg}).
\item \emph{Exhausts a box}: some node-scoped \bbox{} has a nonempty node scope contained in the
  set of removed nodes.
\item \emph{Link boxes of a port}: the chain \bboxes{} (\cref{ch:diagrams}) whose link list
  contains the port.
\end{itemize}

\paragraph{Anchors.}
An anchored search restricts candidates to those whose seed nodes meet the $h$-hop closure of the
anchor set, where $h$ is $0$ for the matchers seeded at a wire and $1$ or $3$ for those that read
further. The closure of the anchors is computed by breadth-first search over the undirected
node adjacency and memoized for the last (adjacency map, $h$) pair.

\paragraph{Per-wire-set indexes.}
The port-to-wire map and the per-port claim count are built once per wire set: each is memoized
for the last frozen wire set object it was built from, keyed by object identity.

\subsection{Spider fusion}\label{sec:rewriting-fusion}

The fusion pattern locates two spiders of one color joined by a wire whose two ports agree on
dimension (the spider rule of \cite{coecke2011interacting} and its qudit form \cite{poor2023zxw}).
A pair joined by $k$ wires yields up to one match per wire, each decided alone; applying one
fuses across its own wire and leaves the other $k-1$ wires as self-loops on the merged node.

\begin{figure}[ht]
\centering
\begin{tikzpicture}[zxd]
  \node[zs] (a) at (0,0) {$\alpha$};
  \node[zs] (b) at (0,1.4) {$\beta$};
  \draw[wire] (a) -- (b);
  \draw[wire] (a) -- (-0.5,-0.8); \draw[wire] (a) -- (0.5,-0.8);
  \draw[wire] (b) -- (-0.5,2.2); \draw[wire] (b) -- (0.5,2.2);
  \node at (1.4,0.7) {$=$};
  \node[zs] (c) at (2.6,0.7) {\tiny$\alpha{+}\beta$};
  \draw[wire] (c) -- (2.1,-0.1); \draw[wire] (c) -- (3.1,-0.1);
  \draw[wire] (c) -- (2.1,1.5); \draw[wire] (c) -- (3.1,1.5);
\end{tikzpicture}
\caption{Spider fusion. The consumed wire joins the two spiders; every other port survives on
the merged spider.}\label{fig:rewriting-fusion}
\end{figure}

\begin{definition}[Fusion match]\label{def:rewriting-fusionmatch}
A \py!FusionMatch(a_id, b_id, wire, shared_dim, side_condition_outcomes, dimension_constraints=(), bindings={})!
records one occurrence. \py!a_id! is always the lower node id and \py!b_id! the higher.
\py!shared_dim! is the settled dimension of the merged spider; \py!bindings! is the accumulator of
every concrete symbol binding the resolution produced. The hash folds \py!bindings! as an unordered
set of items, so equal matches hash equally within one process. Its support is the two node ids,
ascending.
\end{definition}

The fusion pattern declares nine side conditions, in this order (\py!FUSION_SIDE_CONDITIONS!):
\begin{enumerate}
\item \code!distinct_nodes!: the endpoints are distinct nodes. Always true; a recorded fact.
\item \code!same_generator_type!: both nodes carry the identical registered generator type, named
  \code!Z! or \code!X!.
\item \code!parallel_wires_become_self_loops!: always true; records the number of other wires
  joining the pair.
\item \code!consumed_wire_direction_permitted_for_color!: for X the consumed wire runs from an
  output port to an input port; for Z any direction combination is permitted.
\item \code!consumed_ports_singly_claimed!: neither consumed port carries a second wire or appears
  on a boundary list.
\item \code!bang_box_scope_agreement!: both nodes have the same innermost node-scoped box (both
  \py!None!, or the same id).
\item \code!dimension_agreement!: the two consumed ports, then every surviving port of both nodes,
  unify against the running shared dimension. A deferred result or a success that binds a symbol
  is accepted and recorded as a constraint; a failure is a non-match.
\item \code!phase_dimension_agreement!: every present phase vector unifies with the shared
  dimension. A deferred result is rejected, and so is a binding to a non-concrete dimension.
\item \code!dimension_guards_satisfied!: every guard of the rule passes at the settled shared
  dimension. An undecided guard blocks the match, with \py!deferred! false.
\end{enumerate}

\begin{signature}
\py!resolve_fusion_match(diagram, a_id, b_id, wire, *, dimension_guards=(), index=None) -> FusionResolution!
\end{signature}
This is the single verification predicate of fusion: the matcher calls it once per candidate,
and the fusion builder (\cref{ch:rules}) calls it again on the working diagram and builds only
from its result. It returns a frozen
\py!FusionResolution(passed, shared_dim, bindings, dimension_constraints, outcomes)!. The
outcomes always cover the nine condition names in declared order; a condition that an earlier
failure prevented from running is recorded as failing, with the detail
\code!not evaluated: <condition> failed first!. When \py!passed! is false, \py!shared_dim! is
\py!None! and the other two fields are empty.

It raises \py!RewriteGrammarError! when \py!a_id == b_id!, when either node is absent, when the
wire's endpoint nodes are not exactly $\{a, b\}$, when an endpoint names an out-of-range port, or
when the wire is not an element of \py!diagram.wires! (a freestanding wire between two real ports
is rejected).

\begin{algorithm}[Fusion resolution]\label{alg:rewriting-fusion}
\leavevmode
\begin{enumerate}
\item Validate the request as above. Let $r_a$ be the wire endpoint on $a$ and $r_b$ the one on
  $b$; count the other wires joining $\{a,b\}$.
\item Record conditions 1 and 3 as passing. Decide condition 2; on failure fail the rest.
\item Condition 4: the consumed ports have the same direction only if the color is Z. On failure
  fail the rest.
\item Condition 5: each consumed port has no other wire claim and no boundary entry. On failure
  fail the rest.
\item Condition 6: compare the innermost node-scoped boxes. On failure fail the rest.
\item Seed the shared dimension $s$ with the dimension of $r_a$ (the lower-id node) and the
  binding map $\beta$ with $\emptyset$.
\item If the diagram's parameter environment (\py!Diagram.resolve_dim!) changes either consumed
  dimension and the resolved pair fails to unify, fail conditions 7 and 8 as a refuted equality.
\item Repeat at most $32$ passes (\py!_MAX_FIXPOINT_PASSES!):
  \begin{enumerate}
  \item Unify the two consumed dimensions, each resolved through the concrete entries of $\beta$,
    against each other (source \code!CONNECTING_PAIR!).
  \item For node $a$, then node $b$, for every port other than the consumed one, inputs before
    outputs and in index order: resolve its dimension through $\beta$ and unify it with $s$
    (source \code!SURVIVING_LEG!).
  \item Snapshot $(s,\beta)$ as the port-verified state.
  \item For the phase of $a$, then of $b$, if present: resolve the phase dimension and unify it
    with $s$ (source \code!NODE_PHASE!). A deferred result, a failure, or a binding to a
    non-concrete dimension fails the resolution: condition 7 is reported passing from the
    port-verified snapshot, condition 8 failing with a cause-specific detail, condition 9 failing.
  \item After every successful unification that binds: merge the concrete bindings into $\beta$
    (a rebinding of a symbol to a different concrete value fails the candidate) and substitute
    them into $s$.
  \item Stop when a full pass changes neither $s$ nor $\beta$.
  \end{enumerate}
  Any unification failure in (a) or (b) fails conditions 7 and 8 with a detail naming the side.
  When the budget runs out, both fail with a budget-exhaustion detail.
\item Re-verify from scratch that both consumed dimensions, every surviving port, and every
  phase, resolved under the final $\beta$, unify with the final $s$ without failure.
\item Record condition 7 as passing, with \py!deferred! true iff the final record holds a
  deferred constraint from the connecting pair or a surviving port.
\item Trial-reattach each present phase to $s$ under $\beta$ (\py!reattach_phase!). If an entry
  index falls out of range, fail condition 8.
\item Record condition 8 as passing, listing the bindings that phase checks made.
\item Evaluate the guards (condition 9). Return \py!passed! true iff it passes.
\end{enumerate}
\end{algorithm}

The fixpoint stops only when both $s$ and $\beta$ are stable: stopping on $s$ alone lets a binding
that touches no symbol of $s$ escape re-checking of earlier ports. Termination follows because
$\beta$ only grows and draws its keys from the finite set of free symbols of the two nodes.

\paragraph{The constraint record.}
During resolution the constraints live in a record keyed by constraint source, in first-derivation
order. A re-check of a source overwrites its entry in place. Its policy over (previous entry, this
check's result) is:
\begin{center}\small
\begin{tabular}{@{}lll@{}}
\toprule
Previous entry & This check & Effect \\
\midrule
none & deferred or bound & record it \\
none & bare identity & nothing \\
\code!DEFERRED! & deferred or bound & overwrite at the current operands \\
\code!DEFERRED! & bare identity & delete the entry \\
\code!BOUND! & deferred or bound & overwrite \\
\code!BOUND! & bare identity & keep the bound entry \\
\bottomrule
\end{tabular}
\end{center}
A bound entry survives a later identity because the identity holds only under that binding. For
a surviving port or the connecting pair, \py!assumed! is the resolved port dimension; for a
surviving port and a phase, \py!equal_to! is the shared dimension as of the check; for the
connecting pair it is the other resolved consumed dimension.

\begin{signature}
\py!reattach_phase(phase: PhaseVector, shared_dim: Dim, bindings: Mapping[str, Dim]) -> tuple[PhaseVector, Mapping[str, Dim]]!
\end{signature}
substitutes the concrete entries of \py!bindings! (as integers) into the phase's entries, then
builds a new phase vector over \py!shared_dim! with the substituted entries, rule, weights, and
weight rule. It also returns the subset of \py!bindings! actually used: the names free in some
entry, extended by the phase's own free symbols when the phase carries a rule or weights. It
raises \py!PhaseDomainError! when an entry index falls outside \py!shared_dim! (\cref{ch:algebra}).

\begin{signature}
\py!find_matches(diagram, *, dimension_guards=(), anchors=None) -> tuple[FusionMatch, ...]!
\end{signature}
\begin{algorithm}[Fusion matcher]\label{alg:rewriting-findfusion}
\leavevmode
\begin{enumerate}
\item Take the wires in sorted order. Validate both endpoints of every wire, then every boundary
  entry (inputs first), raising \py!RewriteGrammarError! for an absent node (\code!references node id .. absent from the diagram!) or an out-of-range index (\code!.. is out of range for node ..!).
\item Group the non-self-loop wires by their unordered endpoint pair.
\item For each candidate wire: order its endpoints as (lower, higher); skip it if anchors are
  given and neither endpoint is an anchor; skip it unless both nodes carry one registered Z or X
  type; call \py!resolve_fusion_match! and keep the candidate iff it passes.
\item Sort by (lower id, higher id, direction and index of the consumed port on the lower node,
  direction and index on the higher node), directions compared by their string values.
\end{enumerate}
\end{algorithm}

\py!FusionPattern(dimension_guards=())! has locality radius $1$ and orders matches by the key of
step~4. \py!innermost_node_scope_box(diagram, node_id) -> BangBoxId | None! is public.


\paragraph{Detail strings.}
The \py!detail! of each passing fusion outcome is part of the certificate text. With $a$, $b$ the
node ids and $r_a$, $r_b$ the consumed ports printed as \py!PortRef! \py!repr!s:
\begin{center}\small
\begin{tabularx}{\textwidth}{@{}>{\raggedright\arraybackslash}p{0.34\textwidth}X@{}}
\toprule
Condition & Passing detail \\
\midrule
\texttt{distinct\_\allowbreak{}nodes} & \code!<a> != <b>! \\
\texttt{same\_\allowbreak{}generator\_\allowbreak{}type} & \code!both nodes are 'Z'! (or \code!'X'!) \\
\texttt{parallel\_\allowbreak{}wires\_\allowbreak{}become\_\allowbreak{}self\_\allowbreak{}loops} & \code!N other wire(s) join the two nodes, surviving as self-loop(s) on the merged spider! \\
\texttt{consumed\_\allowbreak{}wire\_\allowbreak{}direction\_\allowbreak{}permitted\_\allowbreak{}for\_\allowbreak{}color} & \code!<r_a> (direction=<dir>) -> <r_b> (direction=<dir>)!, followed for a same-direction Z wire by a note that it is permitted for Z only \\
\texttt{consumed\_\allowbreak{}ports\_\allowbreak{}singly\_\allowbreak{}claimed} & \code!neither <r_a> nor <r_b> is claimed by another wire or listed on a boundary! \\
\texttt{bang\_\allowbreak{}box\_\allowbreak{}scope\_\allowbreak{}agreement} & \code!neither node is inside a bang box! or \code!both nodes are inside bang box <id>! \\
\texttt{dimension\_\allowbreak{}agreement} & \code!<x> == <y>! for an identity; \code!<x> == <y> (deferred, assumed)!; \code!<x> == <y> (bound: s := v, ...)!; or a note that the binding is non-concrete; followed, when surviving ports recorded constraints, by \code!; surviving leg(s) resolved to shared_dim=<s> with N additional assumed dimension equality/ies! \\
\texttt{phase\_\allowbreak{}dimension\_\allowbreak{}agreement} & \code!no phase present on either node!, or \code!present phase dimension(s) unify with the resolved shared leg dimension <s>!, followed by \code!; assuming s := v, ...! when phase checks bound symbols \\
\texttt{dimension\_\allowbreak{}guards\_\allowbreak{}satisfied} & \code!no dimension guards declared!, or \code!dimension guard(s) <g, ...> hold at the matched shared dimension <s>! \\
\bottomrule
\end{tabularx}
\end{center}
A failing guard outcome reads \code!dimension guard(s) <g, ...> do not hold at the matched shared dimension <s>!;
an undecided one reads \code!dimension guard(s) <g, ...> cannot be decided at the non-concrete matched shared dimension <s>, so the rewrite is blocked rather than assumed!.
The connecting-pair failure reads \code!<x> != <y>: the connecting pair does not unify!, and a
surviving-port failure \code!a surviving leg of the A-side node (<x>) does not unify with shared_dim (<s>)! (B-side likewise).

\begin{example}[Fusion at three dimension pairs]\label{ex:rewriting-fusion}
Let node $0$ be a Z spider with input dimension $d$ and output dimension $d_a$, and node $1$ a Z
spider with input dimension $d_b$ and output dimension $d$, the output of $0$ wired to the input of
$1$, the remaining two ports on the boundary.
\begin{itemize}
\item $d_a = d_b = d$: one match with \py!shared_dim! $d$, no bindings, no constraints; the
  \code!dimension_agreement! detail is \code!d == d!.
\item $d_a = d$, $d_b = e$: unification binds $d := e$, a non-concrete binding, so it is recorded
  as the constraint \code!d == e (bound, connecting pair)! with equation \code!d = e!, is not added
  to the binding map, and leaves \py!shared_dim! at the seed $d$. The detail notes that the binding is
  non-concrete and left unused.
\item $d_a = d$, $d_b = 3$: unification binds $d := 3$; the binding enters the map, the shared
  dimension becomes $3$, every surviving port resolves to $3$, and the match carries
  \py!bindings = {"d": 3}! and the constraint \code!d == 3 (bound, connecting pair)!.
\end{itemize}
Applying \code!spider_fusion! in the first case consumes nodes $0$ and $1$, creates node $2$, and
maps input $0$ of node $0$ and output $0$ of node $1$ to input $0$ and output $0$ of node $2$; the
introduced scalar is $1$.
\end{example}

\subsection{Fourier cancellation}

\py!FOURIER_CHAIN_LENGTH! $=4$: four Fourier boxes in series cancel, since $F^4 = 1$. The pattern
declares \code!chain_is_series! and \code!same_dimension!.

\begin{algorithm}[$F^4$ matcher]\label{alg:rewriting-f4}
For every node id in ascending order whose node is a registered Fourier box:
\begin{enumerate}
\item Follow the wire on output port $0$; its far endpoint must be input port $0$ of another
  Fourier box. Append that box unless it is already in the chain. Repeat until the chain holds
  four boxes or the walk stops.
\item Discard chains shorter than four.
\item Take the dimensions of output port $0$ and input port $0$ of each box in chain order. Every
  dimension that differs syntactically from the first must unify with it; a failure discards the
  chain, any other non-identity result records a \code!DEFERRED! constraint with source
  \code!CONNECTING_PAIR!.
\item Discard the chain when anchors are given and are disjoint from it, or when its four nodes do
  not share one set of node boxes.
\item Emit \py!FourierMatch(node_ids, wires, shared_dim, outcomes, constraints)!; the
  \code!same_dimension! outcome is deferred iff a constraint was recorded.
\end{enumerate}
\end{algorithm}
The matches are ordered by the first node id. A chain of five boxes yields two overlapping
matches. \py!FourierCancellationPattern! has locality radius $3$.

\subsection{Cap}

\py!find_cap_matches(diagram, *, anchors=None, swapped=False)! locates a phaseless Z state (no
input, one output) wired into a phaseless X effect (one input, no output), or X into Z when
\py!swapped!. The function \py!colour_pair(swapped)! returns \code!("Z", "X")!, or
\code!("X", "Z")! when swapped; every swapped matcher in this section reads its two spider roles
from it. A wire qualifies when its output-side endpoint is the state, both nodes are distinct and
registered with the required names and arities, neither carries any phase vector (here a phase
vector of zeros also disqualifies), each joined port has exactly one wire claim and no boundary
entry, the two port dimensions are syntactically equal, and both nodes lie outside every \bbox{}.
Matches are ordered by (state id, effect id); \py!CapPattern(swapped=False)! has radius $1$. Its
conditions are \code!state_is_a_z_spider!, \code!effect_is_an_x_spider!,
\code!joined_and_unclaimed!, \code!phaseless!, \code!same_dimension!, and
\code!outside_every_bang_box! (with the colors exchanged in the first two names when swapped).

\subsection{Identity removal}\label{sec:rewriting-identity}

\py!find_identity_matches(diagram, *, anchors=None)! locates a phaseless spider with one input and
one output and splices it out. Its conditions are \code!node_is_a_phaseless_spider!,
\code!one_in_one_out!, \code!not_a_self_loop!, \code!legs_singly_claimed!,
\code!splice_keeps_two_ports!, \code!same_dimension!, \code!leaves_every_bang_box_populated!,
\code!no_port_boxed_leg!, and \code!splice_keeps_copies!; a match records all of them as passing
with the detail \code!re-derived from the diagram!.

\begin{algorithm}[Identity matcher]\label{alg:rewriting-identity}
For every node in ascending id order:
\begin{enumerate}
\item Require a registered phaseless spider (anchored: within one hop of an anchor), with one
  input and one output and a uniform port dimension $s$.
\item Require both ports claimed at most once. Let $w_i$, $w_o$ be the wires on the input and the
  output port. Skip a self-loop ($w_o = w_i$).
\item Choose the spliced wire: $w_o$ when it exists and either $w_i$ is absent or the output port
  is no chain link port; otherwise $w_i$; skip if neither exists. The other port is the
  \emph{surviving} port; the far endpoint of the spliced wire is the \emph{far} port.
\item Skip when the far port lies on the node itself, when anchors are given and meet neither the
  node nor the far node, when the surviving port is a boundary slot and the far port's direction
  differs from the surviving port's, when the surviving port's own neighbor is the far port, or
  when a port of the node lies in a port-scoped \bbox{}.
\item Compute the spliced equation: the dimensions of the surviving and far ports must unify. When
  they are equal, or the surviving port's neighbor already has the surviving dimension, no
  constraint is recorded; otherwise a deferred result records a \code!DEFERRED! and a binding a
  \code!BOUND! constraint (source \code!CONNECTING_PAIR!).
\item When the surviving port has no wire, try to lift the node's box (below). Without a lifted
  box, skip if removing the node exhausts a node-scoped box, and skip unless the splice keeps
  copies paired (below).
\end{enumerate}
\end{algorithm}

\paragraph{Box lifting.}
The innermost box $b$ of the node \emph{lifts} to a port-scoped \bbox{} on the far port when $b$
holds only this node, is not extended, has no child box, the surviving port is a boundary slot,
the far node is a spider with no port in a port-scoped \bbox{}, the far node's boxes equal the
node's boxes minus $b$, and $b$'s parent equals the far node's innermost box. The match then
carries \py!lifted_box = b!.

\paragraph{Copies stay paired.}
Let $M$ be the node's boxes and, for each wired port, $N$ the far node's boxes minus the link boxes
of the far port. With \py!own_ports! given (identity removal): if both own ports are link ports
the splice is rejected; if exactly one is, the neighbor across the other port must exist, meet it
at no link port, and have boxes exactly $M$. Otherwise the splice is rejected when the two far
ports share a link box, when some $N \not\subseteq M$, or when $M \neq \emptyset$ and no $N$
equals $M$.

\py!IdentityMatch(node_id, wire, surviving_ref, far_ref, shared_dim, side_condition_outcomes, dimension_constraints=(), lifted_box=None)!
has support \{node, both wire endpoints, far node\}. \py!IdentityRemovalPattern! has radius $1$
and orders by node id.

\subsection{Triangle--inverse cancellation}

\py!find_triangle_matches(diagram, *, anchors=None)! locates a triangle $T$ and an inverse
triangle $T^{-1}$ (\cref{ch:diagrams}) in series, either order, joined from output port $0$ of the
first to input port $0$ of the second. Both nodes must be registered, one-in-one-out, with the
joined ports claimed exactly once by the wire, every port of both nodes of one dimension, and both
outer ports (input of the first, output of the second) claimed at most once. The spliced wire is
the second node's output wire when present, otherwise the first node's input wire; neither
present skips. The far port may not lie on either node, the surviving port's neighbor may not be
the far port, both nodes must have one innermost box, the pair may not exhaust a node-scoped box,
the splice must keep copies paired (without \py!own_ports!), and the spliced equation must have a
solution. Matches are ordered by (first id, second id); the pattern has radius $1$. Its conditions
are \code!inverse_pair_in_series!, \code!joining_wire_unclaimed!, \code!outer_legs_singly_claimed!,
\code!splice_keeps_two_ports!, \code!same_dimension!, \code!bang_box_scope_agreement!, and
\code!leaves_every_bang_box_populated!.

\subsection{State copy}

\py!find_state_copy_matches(diagram, *, anchors=None, swapped=False)! locates a phaseless X state
(no input, one output) wired into input port $0$ of a phaseless Z spider with exactly one input
(any number of outputs), or Z into X when swapped. Both joined ports must be claimed exactly once
by the wire; the state's output and all ports of the spider share one dimension; both nodes lie
outside every \bbox{}. \py!StateCopyMatch! records \py!output_count!, the spider's number of
outputs. Ordered by (state id, spider id); radius $1$.

\subsection{Hopf law}

\py!find_hopf_matches(diagram, *, anchors=None, swapped=False)! locates a Z spider and an X spider
(X and Z when swapped) joined by a direct wire and by a second path through two Fourier boxes in
series, the qudit Hopf law of~\cite{poor2023zxw}.
\begin{algorithm}[Hopf matcher]\label{alg:rewriting-hopf}
For every node $z$ in ascending id order of the upper color, within three hops of an anchor, with a
uniform port dimension $s$:
\begin{enumerate}
\item For each output index $i$ of $z$ whose port is claimed exactly once by a wire: the far port
  must be an input port, claimed exactly once, of a registered node $x$ of the lower color with
  uniform dimension $s$.
\item For each other output index $j$ claimed exactly once: walk two Fourier boxes (each entered
  at input $0$, one-in-one-out, distinct, with entry and exit claimed exactly once), then the wire
  from the second box's output to a far port.
\item Require the far port to be an input port of $x$ different from the direct one and claimed
  exactly once; both boxes of uniform dimension $s$; any phase on $z$ or $x$ over $s$; all four
  nodes outside every \bbox{}; and, if anchors are given, one of the four an anchor.
\end{enumerate}
\end{algorithm}
\py!HopfMatch! carries \py!z_id!, \py!x_id!, \py!fourier_ids!, the four wires, \py!z_leg_indices = (i, j)!,
and \py!x_leg_indices!. Matches are ordered by (z id, x id, z indices, x indices); radius $3$.

\subsection{Bialgebra}

\py!find_bialgebra_matches(diagram, *, anchors=None, swapped=False)! locates a phaseless X spider
with two inputs and one output whose output port $0$ feeds input port $0$ of a phaseless Z spider
with one input and two outputs (Z into X when swapped). Both joined ports are claimed exactly once
by the wire; both nodes have one uniform port dimension, the same for both; both lie outside every
\bbox{}. Ordered by (two-to-one id, one-to-two id); radius $1$.

\subsection{Fourier box on a state}

\py!find_fourier_state_matches(diagram, *, anchors=None, swapped=False)! locates a one-in-one-out
Fourier box in series with a phaseless Z spider (X when swapped) that is a state (no inputs, one
output wired to the box's input $0$) or an effect (one input wired from the box's output $0$).
The spider's port index is $0$; the joined ports are claimed exactly once by the wire; the box's
free port is claimed at most once; all ports of both nodes share one dimension; both nodes have one
innermost box. \py!FourierStateMatch! records \py!free_ref!, the box's other port, and
\py!is_state!. Ordered by (spider id, box id); radius $1$.

%======================================================================
\section{Pattern matching against diagrams}\label{sec:rewriting-template}

The module \module{archytaszx.rewrite.match_template} states a rule as an equation between two
small diagrams, matches either side as a subgraph of a host diagram, and rebuilds the other side in
its place. Most rules of the complete axiom set (\cref{ch:completeness}) are written this way.
In code the small diagram is called a \emph{template}; this manual calls it the pattern
diagram of a side.

\subsection{Pattern diagrams}

\begin{definition}[Template port, node, and diagram]\label{def:rewriting-template}
A template port \py!TPort = tuple[str, Direction, int]! is (node key, direction, index). A
template node is a frozen record
\begin{signature}
\py!TNode(key: str, gen: GeneratorType, ins: int, outs: int, param: Param = PHASELESS, open: bool = False, pink: bool = False, dim: Dim | None = None, port_dims: tuple[tuple[Dim, ...], tuple[Dim, ...]] | None = None)!
\end{signature}
with \py!ins! inputs and \py!outs! outputs. An \py!open! node lets its host node carry further
ports in its symmetric groups (its \term{rest}); an open node has no terminal port. A \py!pink!
node is an X spider in the normalization of~\cite{poor2023zxw}, which equals
$d^{(m+n)/2-1}$ times the engine's X spider of $m$ inputs and $n$ outputs. \py!dim! gives every
port of the node a dimension other than the template's; \py!port_dims! gives the input and
output dimensions of a dimension binder B or splitter S port by port.
\py!leg_dims(default)! returns \py!port_dims! when set, else \py!(dim,)*ins, (dim,)*outs! with
\py!dim! defaulting to \py!default!.

A \py!Template(nodes, wires=(), terminals={}, bare=(), rest={}, scalar=None)! is a frozen
record: template nodes, wires between template ports, named terminals (the ports the interface of
the equation passes through), \py!bare! pairs of terminal names joined directly by a wire (on a
built side only), \py!rest! mapping each open node of the \emph{other} side to the node of this
side that receives its rest, and \py!scalar(dim, bindings) -> Scalar!, a factor this side carries.
\py!Template.node(key)! returns the node with that key or raises \py!RewriteGrammarError!
(\code!template has no node 'k'!). The first node is the \term{anchor}.
\end{definition}

\begin{definition}[Side and axiom]\label{def:rewriting-axiom}
A \py!Side(template, arities=..., min_size=...)! is one side of an equation:
\begin{itemize}
\item \py!template(dim: Dim, arities: Mapping[str, int]) -> Template | None!, the pattern
  diagram at a dimension and arity parameters (such as \code!{"n": 3}!);
\item \py!arities(diagram, node_id, dim) -> Iterable[Arities]!, the arity assignments a host node
  can anchor; default \py!({},)!;
\item \py!min_size(dim) -> int!, a lower bound on the pattern's node count plus wire count;
  default $0$.
\end{itemize}
At construction the \py!template! function is wrapped in a memo keeping its last $256$
(\py!_TEMPLATE_MEMO!) results keyed by (dimension, sorted arity items), least recently used
evicted first; an unhashable key bypasses the memo; wrapping is idempotent.

A \py!TemplateAxiom! is the frozen record of an equation $\mathit{left} = \mathit{right}\cdot s$
applied in one direction:
\begin{signature}
\py!TemplateAxiom(name, left: Side, right: Side, scalar=lambda d, b: Scalar.one(), conditions: tuple[Condition, ...] = (), concrete_only: bool = False, description: str = "", quantifiers: tuple[str, ...] = (), free={}, converse_free={}, solve=None, candidates=None)!
\end{signature}
\py!scalar(dim, bindings)! is the factor the equation states beyond the sides' own scalars and pink
normalizations. A \py!Condition(name, description, holds)! is a named predicate
\py!holds(dim, bindings) -> bool! on the bindings of the matched side. \py!concrete_only!
restricts matching to concrete dimensions. \py!free! maps each variable the built side reads but
the matched side does not bind to a default \py!(dim, bindings) -> value!; a user-supplied value
(the \azx!apply ... with NAME=VALUE! form of \cref{ch:shell}) replaces the default, and the
conditions then check it. \py!solve(dim, bindings)! returns the free variables the conditions
determine from the others, or \py!None!. \py!candidates(dim, bindings)! returns further value
assignments of the free variables, which exhaustive search enumerates.
\end{definition}

\py!axiom.side_conditions! is \py!TEMPLATE_SIDE_CONDITIONS! followed by one \py!SideCondition!
per condition. The five template conditions are \code!template_matches!, \code!parameters_match!,
\code!same_dimension!, \code!outside_every_bang_box!, and \code!right_side_defined!.

\py!axiom.converse(name)! returns the equation read right to left: sides exchanged, scalar
$s^{-1}$ (a zero scalar raises \py!RewriteDomainError!, \code!a template rule scalar of 0 has no converse!),
the same conditions, \py!concrete_only!, and quantifiers, the description suffixed with
\code! Read right to left.!, and \py!free! set to the original \py!converse_free!. The converse
carries no \py!converse_free!, \py!solve!, or \py!candidates!.

\subsection{Parameters}

A template node's \py!param! is a \py!Param! with three operations:
\py!match(phase, dim, bindings) -> Bindings | None! extends the bindings by matching the host's
phase vector, \py!build(dim, bindings) -> PhaseVector | None! constructs the vector, and
\py!variables()!, \py!binds()! name the variables read and bound. A phase vector \py!None! is the
zero vector throughout. The default \py!match! builds the vector and compares.
\begin{itemize}
\item \py!Fixed(vector)!: \py!vector(dim)!. \py!PHASELESS = Fixed(lambda d: None)!.
\item \py!Var(name)!: on match, binds \py!name! to the host vector (the zero vector over $d$ when
  phaseless), or, when already bound, requires the same vector.
\item \py!Derived(reads, compute)!: matches only when every name of \py!reads! is bound and
  \py!compute(dim, bindings)! equals the host vector; an \py!ArithmeticError!, \py!ValueError!, or
  \py!KeyError! in \py!compute! is a non-match.
\item \py!Kj(name, sign=1, offset=0)!: the vector $K_{\sigma j+o}$ (below). On match it reads
  $j_h$, the index of the host vector; if \py!name! is bound to $j$ it requires
  $j_h \equiv \sigma j + o$; otherwise it binds $j = \sigma(j_h - o)$, reduced modulo $d$ at a
  concrete dimension.
\end{itemize}

\begin{definition}[The vector $K_j$]\label{def:rewriting-kj}
For an integer expression $j$ and dimension $d$, $K_j$ is the phase vector with
$jk/d$ turns at index $k$:
\[
  K_j = \bigl(\tfrac{jk}{d}\ \text{turns}\bigr)_{k=1}^{d-1} .
\]
\py!k_vector(j, dim)! returns \py!None! when $j \equiv 0$. At a concrete $d=r$ it requires an
integer $j$ (else \py!RewriteDomainError!), reduces $j_r = j \bmod r$, and returns the explicit
vector with entry $k$ equal to $(j_r k \bmod r)/r$ turns for $k = 1,\dots,r-1$ (memoized for $512$
pairs $(j_r, r)$). At a symbolic $d$ it returns the rule vector with entry $j\cdot k/d$ turns.
\end{definition}

\py!kj_index(phase, dim)! inverts $K$: it returns $0$ for \py!None! or a zero vector and
\py!None! for a weighted vector. At concrete $d=r$ it reads the turns $t_1$ of entry $1$; when
$t_1$ is rational and $j = r t_1$ is an integer and $K_j$ equals the vector, it returns
$j \bmod r$. At symbolic $d$ it requires a rule vector with no explicit entries, computes
$j = \mathrm{simplify}(t \cdot d / k)$ from the rule's turns $t$, and returns $j$ when $j$ is free of
the index symbol, integer-valued, and $K_j$ equals the vector.

Two parameters are the \emph{same vector} when they are equal as phase vectors, or, at a concrete
dimension $r$, when they agree at every index: all of $1,\dots,r-1$ when either carries a rule or a
weight rule, else the union of their explicit indices. Unweighted entries agree when their turns
are equal, or, not both rational, when their difference simplifies to an integer. Weighted entries
agree when their coefficients are equal, or, both constant, when their exact values in
$\mathbb{Q}(\sqrt2)$ coincide (\cref{ch:algebra}), or when their difference simplifies to $0$. At
a symbolic dimension only structural equality counts.

\subsection{Port symmetry}

Matching is subgraph matching up to the symmetry of each generator's ports. A node's ports fall
into \emph{groups} whose members are interchangeable:
\begin{center}\small
\begin{tabular}{@{}lll@{}}
\toprule
Generator & Groups & Open groups \\
\midrule
Z spider, discard G & \code!legs! (all inputs and outputs) & \code!legs! \\
X spider & \code!in! (inputs), \code!out! (outputs) & \code!in!, \code!out! \\
W node & \code!in!, \code!out! & \code!out! \\
F, T, $T^{-1}$, B, S, others & one group \code!<direction><index>! per port & none \\
\bottomrule
\end{tabular}
\end{center}
A host node is \emph{compatible} with a template node when the generator names agree and, for every
template group, the host group has exactly as many ports, or at least as many when the template
node is open and the group is one of its open groups; and the host has no group the template
lacks, unless the template node is open.

A host node can be \emph{placed} at a template node when it is not yet used, is not boxed (it is
in no node scope and no port of it is in a port scope), is compatible, has the template's port
dimensions (port by port for \py!port_dims!; otherwise every port of dimension \py!dim! or the
template dimension, and any phase over the same dimension; a node without ports must carry a phase
over that dimension), and its phase matches the node's parameter. A \py!RewriteDomainError! during
the parameter match is a non-match.

\subsection{Subgraph matching}

\begin{signature}
\py!match_template(diagram, template, anchor: NodeId, dim: Dim, partner=None, boxed=None) -> Iterator[_State]!
\end{signature}
enumerates every complete assignment of the pattern diagram with its first node placed at
\py!anchor!. A state holds the node assignment (key to host id), the used host nodes, the port
assignment (template port to host port), the used host ports, the matched host wires, and the
bindings. \py!partner! maps each wired host port to the far endpoint of its wire.

\begin{algorithm}[Subgraph matching]\label{alg:rewriting-subgraph}
\leavevmode
\begin{enumerate}
\item Place the anchor; stop if it cannot be placed.
\item \emph{Extend}$(S, R)$ with $R$ the unmatched template wires in declaration order:
  \begin{enumerate}
  \item If the key (sorted node assignment, used ports, assigned template ports) was seen in this
    call, return; record it.
  \item If $R$ is empty, \emph{finish}.
  \item Pick the first wire of $R$ with a placed endpoint $a$; if there is none, place the first
    unplaced template node at every host node in ascending id order and extend each placement
    (a new connected component).
  \item The host candidates of $a$ are its assigned host port, or the unused host ports of the
    group of $a$ on its host node. For each candidate $h$ with a far endpoint $f$ that is not used
    and is not $h$ itself:
    \begin{itemize}
    \item if the other endpoint $b$ is placed, require $f$ to lie on $b$'s host node and to be
      $b$'s assigned port or one of its candidates; assign $a\mapsto h$, $b\mapsto f$, record the
      wire, and extend;
    \item otherwise place $b$'s node at $f$'s node, require $f$ to be a candidate of $b$, assign,
      record, and extend.
    \end{itemize}
  \end{enumerate}
\item \emph{Finish}: if some template node is unplaced, place a new component as in (c) with no
  wires left. Otherwise, for every node that is not open and every group, pair the unassigned
  template ports with the unused host ports of the group; the counts must agree, and the host
  ports are taken in order (node id, direction value, index). Yield the state.
\end{enumerate}
\end{algorithm}

The search is a depth-first backtracking over template wires, the anchor-first strategy of
!-graph matching~\cite{kissinger2014pattern} restricted to fixed arities: a family of patterns
with arity parameters is matched by trying each arity assignment the side enumerates.

\subsection{The template pattern}

\py!TemplatePattern(axiom)! is a \py!Pattern! of locality radius $8$. Its \py!repr! is
\code!TemplatePattern('<name>', concrete_axioms=<CONCRETE_AXIOMS>)!.

\paragraph{The concrete switch and resolved views.}
The module flag \py!CONCRETE_AXIOMS!, default \py!True!, controls two things: when false, a
\py!concrete_only! axiom finds no matches, and matching reads the host diagram as is. When true,
matching reads the \emph{resolved view} of the host: the diagram with every dimension symbol bound
in its parameter environment and occurring in some port or phase dimension substituted
(\py!Diagram.substitute!), node ids preserved; a diagram with no parameters is its own view.
Inside a \py!resolved_view_scope()! context, one view is memoized per diagram object.

\begin{algorithm}[Template matcher]\label{alg:rewriting-templatefind}
\py!find_matches(host)!:
\begin{enumerate}
\item Return no match when the axiom is concrete-only and \py!CONCRETE_AXIOMS! is false. Let $D$
  be the resolved view of the host.
\item Build the left pattern at symbol $d$ with no arities. If that succeeds and has no nodes, the
  side is an \emph{insertion}: return the insertion matches (below).
\item For every node id of $D$ in ascending order as the anchor: let the dimension be that of its
  first port, or of its phase when portless; skip when absent (in $D$ or in the host), when the
  axiom is concrete-only and the dimension is symbolic, or when \py!min_size! exceeds the node count
  plus the wire count of $D$.
\item For every arity assignment of the side at this anchor: build the pattern (once per dimension
  and arities within this call) under a \emph{template budget} equal to the node count plus wire
  count of $D$; a pattern construction that calls \py!check_template_budget(size)! with a larger
  size raises \py!TemplateTooLarge! and counts as no pattern. Skip an empty pattern or one whose
  anchor generator differs from the host node's.
\item For every state of \py!match_template!, convert it to a match (below) and keep the first match
  per key (support set, matched wire set, arities).
\item Return the matches sorted by the order key (support ids, arities, sorted wire keys).
\end{enumerate}
\end{algorithm}


\paragraph{Anchored template search.}
\py!TemplatePattern.find_matches_anchored(diagram, anchors)! returns the matches whose support
meets \py!anchors!. It returns no match under the concrete switch as above. It falls back to the
default (a full scan filtered to the anchors) when the resolved view is a different object than the
host (the diagram has bound parameters), when the left side is an insertion, or when some pattern
built for some anchor and arity assignment is \emph{disconnected}: its wires, read undirected, do
not join all its nodes (checked by a depth-first walk from the first node). Otherwise:
\begin{enumerate}
\item Build the plan of (anchor, dimension, arities, pattern) exactly as steps 3 and 4 of
  \cref{alg:rewriting-templatefind}, and let the reach be the largest node count minus one among the
  planned patterns.
\item Compute, by breadth-first search over the partner map from the anchor set, the wire-hop
  distance of every node at most the reach away.
\item Run \py!match_template! only from a planned anchor whose distance is at most its pattern's
  node count minus one; convert and deduplicate as in the full search, with the host dimension equal
  to the pattern dimension.
\item Return the matches sorted by the order key whose support meets the anchors.
\end{enumerate}
A connected pattern of $n$ nodes spans at most $n-1$ hops from its anchor, so every match whose
support meets the anchors is found from an anchor within that distance.

\begin{definition}[Template match]\label{def:rewriting-templatematch}
\py!TemplateMatch(rule_name, arities, nodes, terminals, rest, wires, bindings, shared_dim, side_condition_outcomes, value_dim=None, dimension_constraints=())!
is a frozen record. \py!arities!, \py!nodes! (key, host id), \py!terminals! (name, host port), and
\py!bindings! (name, value) are tuples sorted by their first element; \py!rest! lists, per open
node, its host ports not matched by the pattern; \py!wires! are the matched host wires, sorted.
\py!shared_dim! is the host dimension of the anchor; \py!value_dim! is the dimension the pattern was
read at (in the resolved view) when it differs, else \py!None!; \py!template_dim! returns
\py!value_dim! or \py!shared_dim!. The support is the matched node ids, or, for an insertion on a
wire, the terminal nodes. \py!binding_map()! returns the bindings as a dictionary.
\end{definition}

\paragraph{Conversion of a state.}
A state becomes a match when: no matched node has a wire to a boxed node (checked only when the
diagram has boxed nodes); every \py!free! default evaluates; the right pattern exists; every bare
wire of the right side has a solvable spliced equation (the splice chooses the far port of the
first terminal when that lies off the match and is not the second terminal, else the far port of
the second, as in the identity matcher); every right-side parameter and scalar evaluates; the left
side's own scalar is not zero; and every condition holds. An \py!ArithmeticError!,
\py!ValueError!, \py!KeyError!, or \py!RewriteDomainError! in any of these is a non-match. The
spliced equations become the match's dimension constraints. All outcomes pass.

\paragraph{Insertion matches.}
A left side with no nodes inserts its right side.
\begin{itemize}
\item With a bare pair $(t_1,t_2)$: one match per unboxed wire (in sorted wire order) and arity
  assignment, at the dimension of the wire's first endpoint, with terminals $t_1\mapsto$ first
  endpoint, $t_2\mapsto$ second. Then, when the axiom has no arity parameters (its right side
  builds at $d$ without arities), one match per boundary port, in order (node id, direction,
  index), whose node is unboxed and not mixed-only (\cref{ch:mixed}); its \py!nodes! is
  \py!(("_copy", owner),)! and both terminals map to the boundary port.
\item Without a bare pair: one match per distinct host dimension, taken from nodes in ascending id
  order.
\end{itemize}

\paragraph{Verification.}
\py!verify(diagram, match)! is true when the match is rediscovered. Otherwise, for an axiom with
free variables, it looks for a rediscovered match with the same nodes, terminals, rest, wires, and
arities whose bound (non-free) values agree with the given match and which equals it after
replacing its bindings; it then checks that every condition holds, every built parameter and
scalar evaluates, and the right pattern exists.

\py!with_values(diagram, match, values)! sets free variables: it raises \py!RewriteGrammarError!
for a name that is not free (\code!<rule>: [names] are not free variables of this rule (free: ..)!)
and \py!RewriteDomainError! (\code!<rule>: the given values do not satisfy the rule's conditions!)
when the result does not verify. \py!candidate_matches(diagram, match)! returns the match with each
\py!candidates! assignment that verifies, distinct and different from the input.
\py!algebraic_matches(diagram, match, values)! applies at a concrete dimension $r$ only: for every
free variable whose default is a phase vector, it enumerates every weight vector with entries in
\py!values! at indices $1,\dots,r-1$ (all $|\mathit{values}|^{r-1}$ of them); with \py!solve! set
and more than one such variable, it enumerates the first variable and asks \py!solve! for the
rest, enumerating the rest only where \py!solve! returns \py!None! or raises.

\subsection{Scalar and rebuild}

\begin{signature}
\py!template_scalar(axiom, match, diagram=None) -> Scalar!
\end{signature}
returns
\begin{equation}\label{eq:rewriting-template-scalar}
  s \;=\; s_{\mathrm{ax}}(d,\beta)\cdot
  \frac{\prod_{v\in R\text{ pink}} D_v^{\,\ell_v/2-1}}{\prod_{v\in L\text{ pink}} D_v^{\,\ell_v/2-1}}
  \cdot \frac{s_R(d,\beta)}{s_L(d,\beta)} ,
\end{equation}
where $d$ is the pattern dimension, $\beta$ the bindings, $D_v$ the node's own dimension or the
host dimension, and $\ell_v$ the port count of node $v$: on the left its pattern ports plus its
rest, on the right its pattern ports plus the rest of every left node it receives. A missing side
scalar is $1$. A non-\py!TemplateMatch! raises \py!RewriteGrammarError!.

\py!template_builder(axiom)! returns a builder carrying \py!side_conditions = axiom.side_conditions!.
It raises \py!RewriteGrammarError! for a match of another rule, runs the coverage check, raises
\py!RewriteDomainError! (\code!<rule>_builder: the match is not among those rediscovered in this diagram!)
when the match does not verify, and builds.

\begin{algorithm}[Rebuild of the right side]\label{alg:rewriting-rebuild}
Let $d_v$ be the pattern dimension, $d$ the host dimension, $\beta$ the bindings.
\begin{enumerate}
\item A match whose first node key is \code!_copy! is a boundary insertion (below).
\item For each right node in order: list its ports, inputs then outputs. A Z or discard port that
  is a terminal mapped to a host boundary entry takes that entry's direction. Append the rest ports
  of every left node this node receives, keeping each port's direction (a W node takes them as
  outputs). Build the phase from the parameter at the node's dimension (or $d_v$) and re-express it
  over $d$ unless the node has its own dimension; a portless Z or X spider without a phase gets the
  zero vector. Add the node with its port dimensions; map each rest port to its new port.
\item For each right terminal: when the host terminal port lies off the match, record a wire from
  it to the new port. When it lies on a consumed node and is a boundary entry whose direction
  differs from the new port's, insert a one-in-one-out Z spider as an adapter: the boundary maps to
  the adapter's port of the boundary direction and the adapter's other port is wired to the new
  port. Otherwise map the host port to the new port.
\item For each bare pair $(t_1, t_2)$ with host ports $h_1, h_2$ and far ports $f_1, f_2$:
  if $f_1$ is off the match and differs from $h_2$, consume the wire $(h_1,f_1)$ and map $h_2$ onto
  $f_1$; else if $f_2$ qualifies symmetrically, consume $(h_2, f_2)$ and map $h_1$ onto $f_2$; a
  mapping onto a far port of the other direction than a boundary entry's goes through an adapter
  Z spider. Else, when both $h_1, h_2$ are boundary entries, add a phaseless Z spider with one port
  per entry, in the entries' directions. Else raise \py!RewriteDomainError!
  (\code!a bare wire of the right-hand side closes on the match itself!).
\item New wires are the right pattern's wires, then the adapter and off-match terminal wires.
\item Rewrite every port-scoped \bbox{} containing a mapped port to contain the mapped image.
\item Return the build result: consumed nodes (sorted), the matched wires plus the spliced ones,
  the port mapping, the scalar of \cref{eq:rewriting-template-scalar}.
\end{enumerate}
\end{algorithm}


\begin{example}[An axiom as two pattern diagrams]\label{ex:rewriting-template}
The axiom D1 of~\cite{poor2023zxw} states that a Z box with parameter $a$ followed by two Fourier
boxes equals two Fourier boxes followed by the Z box with the reversed parameter. As pattern
diagrams, with ports written (key, direction, index):
\begin{lstlisting}[style=python]
left = Template(
    nodes=(TNode("z", Z_SPIDER, 1, 1, Var("a")),
           TNode("f0", FOURIER_BOX, 1, 1), TNode("f1", FOURIER_BOX, 1, 1)),
    wires=((("z", OUT, 0), ("f0", IN, 0)), (("f0", OUT, 0), ("f1", IN, 0))),
    terminals={"in": ("z", IN, 0), "out": ("f1", OUT, 0)})
right = Template(
    nodes=(TNode("f0", FOURIER_BOX, 1, 1), TNode("f1", FOURIER_BOX, 1, 1),
           TNode("z", Z_SPIDER, 1, 1,
                 Derived(("a",), lambda d, b: reversed_vector(b["a"], d)))),
    wires=((("f0", OUT, 0), ("f1", IN, 0)), (("f1", OUT, 0), ("z", IN, 0))),
    terminals={"in": ("f0", IN, 0), "out": ("z", OUT, 0)})
ZXW_D1 = TemplateAxiom("zxw_d1", Side(lambda d, a: left), Side(lambda d, a: right),
                       concrete_only=True)
\end{lstlisting}
Matching anchors \code!z! at every Z spider with one input and one output whose ports carry one
concrete dimension; \py!Var("a")! binds its phase vector; the two wires are then followed through
the partner map to two Fourier boxes. The rebuild adds the two Fourier boxes and the new Z spider,
maps the host port of terminal \code!in! to input $0$ of the new \code!f0! and that of \code!out! to
output $0$ of the new \code!z!, and wires the three in series. The identity S2 has a one-node left
side with terminals \code!in! and \code!out! and a right side with no nodes and
\py!bare=(("in", "out"),)!; its rebuild splices the far port of \code!in! onto that of \code!out!.
\end{example}

\paragraph{Boundary insertion.}
The node owning the boundary port is rebuilt unchanged as a copy; each of its other ports maps to
the copy's. The right pattern's nodes are added at the port's dimension. Of the two terminal ports,
the one whose direction equals the boundary port's faces the boundary (the second terminal is
tried first); the boundary port maps to it, and the copy's port is wired to the other terminal.
When neither faces the boundary, \py!RewriteDomainError! is raised.

\paragraph{Construction helpers.}
\py!legs_of(diagram, node_id)! returns (input count, output count); \py!neighbours(diagram, node_id)!
returns the far node of every wired port in port order, with repetition;
\py!ranges(*bounds)! is the Cartesian product.

%======================================================================
\section{W-node and connective matchers}\label{sec:rewriting-zxw}

The module \module{archytaszx.rewrite.match_zxw} implements the fixed-shape matchers for the W
node, the triangle, and the dimension connectives B and S of qufinite
ZXW~\cite{wang2023qufinitezxw}. Every pattern here requires its ports to carry exactly (syntactically)
equal dimensions: the \emph{shared dimension} of a set of nodes is the single element of the set
of dimensions of all their ports, and a match requires it to exist. A node \emph{is} a generator
when its type is registered and carries that generator's name. A \emph{phaseless state} of a color
is a phaseless spider of that color with no input and one output; a \emph{phaseless effect} has one
input and no output. Each matcher is wrapped in a \py!Pattern! subclass that delegates to it, as in
\cref{sec:rewriting-fixed}; every match records all of its declared conditions as passing.


\begin{figure}[htbp]
\centering
\begin{minipage}[b]{0.23\textwidth}\centering
\begin{tikzpicture}[zxd]
  \node[wnode] (w1) at (0,0.6) {}; \node[wnode] (w2) at (-0.4,1.5) {};
  \draw[wire] (0,0) -- (w1); \draw[wire] (w1) -- (w2); \draw[wire] (w1) -- (0.5,2.3);
  \draw[wire] (w2) -- (-0.8,2.3); \draw[wire] (w2) -- (-0.1,2.3);
\end{tikzpicture}\\(a) W fusion
\end{minipage}\hfill
\begin{minipage}[b]{0.23\textwidth}\centering
\begin{tikzpicture}[zxd]
  \node[xdot] (s) at (0,0.4) {}; \node[wnode] (w) at (0,1.3) {};
  \draw[wire] (s) -- (w); \draw[wire] (w) -- (-0.45,2.2); \draw[wire] (w) -- (0.45,2.2);
\end{tikzpicture}\\(b) X state into W
\end{minipage}\hfill
\begin{minipage}[b]{0.23\textwidth}\centering
\begin{tikzpicture}[zxd]
  \node[tri] (t) at (0,0.9) {}; \node[xdot] (e) at (0,1.9) {};
  \draw[wire] (0,0) -- (t) -- (e);
\end{tikzpicture}\\(c) $T$ into X effect
\end{minipage}\hfill
\begin{minipage}[b]{0.23\textwidth}\centering
\begin{tikzpicture}[zxd]
  \node[snode] (s) at (0,0.6) {S}; \node[bnode] (b) at (0,1.9) {B};
  \draw[wire] (0,0) -- (s);
  \draw[wire] ($(s.north)+(-0.2,0)$) -- ($(b.south)+(-0.2,0)$);
  \draw[wire] ($(s.north)+(0.2,0)$) -- ($(b.south)+(0.2,0)$);
  \draw[wire] (b) -- (0,2.6);
\end{tikzpicture}\\(d) S then B
\end{minipage}

\bigskip
\begin{minipage}[b]{0.23\textwidth}\centering
\begin{tikzpicture}[zxd]
  \node[zdot] (a) at (-0.25,0.4) {}; \node[zdot] (c) at (0.25,0.4) {};
  \node[bnode] (b) at (0,1.4) {B};
  \draw[wire] (a) -- ($(b.south)+(-0.25,0)$); \draw[wire] (c) -- ($(b.south)+(0.25,0)$);
  \draw[wire] (b) -- (0,2.2);
\end{tikzpicture}\\(e) two Z states into B
\end{minipage}\hfill
\begin{minipage}[b]{0.23\textwidth}\centering
\begin{tikzpicture}[zxd]
  \node[zdot] (z) at (0,0.8) {}; \node[wnode] (w) at (0,1.7) {};
  \draw[wire] (-0.45,0) -- (z); \draw[wire] (0.45,0) -- (z);
  \draw[wire] (z) -- (w); \draw[wire] (w) -- (-0.45,2.5); \draw[wire] (w) -- (0.45,2.5);
\end{tikzpicture}\\(f) W/Z bialgebra
\end{minipage}\hfill
\begin{minipage}[b]{0.23\textwidth}\centering
\begin{tikzpicture}[zxd]
  \node[wnode] (w) at (0,0.7) {}; \node[zdot] (e) at (-0.45,1.6) {};
  \draw[wire] (0,0) -- (w); \draw[wire] (w) -- (e); \draw[wire] (w) -- (0.45,2.3);
\end{tikzpicture}\\(g) W output into Z effect
\end{minipage}\hfill
\begin{minipage}[b]{0.23\textwidth}\centering
\begin{tikzpicture}[zxd]
  \node[zs] (z) at (0,1.0) {};
  \draw[wire] (-0.5,0) -- (z); \draw[wire] (0.5,0) -- (z); \draw[wire] (z) -- (0,2.0);
  \draw[bang] (0.25,-0.1) rectangle (0.6,0.55); \node[banglabel] at (0.6,0.55) {$b$};
\end{tikzpicture}\\(h) port-boxed port
\end{minipage}
\caption{Left-hand shapes of the W-node, connective, and port-box matchers. In (h) the dashed
outline marks the single port that forms the scope of a top-level port-scoped \bbox{}.}
\label{fig:rewriting-zxw-shapes}
\end{figure}

The matchers are listed below with their pattern class and locality radius, the shape they locate, their conditions, and the sort key of their result.
\begin{description}[style=nextline,leftmargin=1.5em]
\item[\py!find_w_fusion_matches!, \py!WFusionPattern!, 1.] Output port $p$ of a W node with one input wired to the input of a second W node with one input;
both joined ports claimed exactly once by the wire; one shared dimension; both nodes with one
innermost node-scoped box; the pair does not exhaust a node-scoped box. \py!WFusionMatch(first_id, second_id, position=p, wire, shared_dim, ...)!. Order (first, second, $p$). Conditions
\code!two_w_nodes!, \code!joining_wire_unclaimed!, \code!same_dimension!,
\code!bang_box_scope_agreement!, \code!leaves_every_bang_box_populated!.
\item[\py!find_w_identity_matches!, \py!WIdentityPattern!, 1.] A W node with one input and one output, uniform dimension, both ports claimed at most once, not a
self-loop. The output wire is spliced when present, else the input wire. The far port is off the
node; the surviving port's neighbor is not the far port; the node does not exhaust a node-scoped
box; no port of the node is port-boxed; the splice keeps copies paired (without \py!own_ports!);
the spliced equation is solvable. Returns \py!IdentityMatch! records without a lifted box. Order:
node id.
\item[\py!find_fed_state_matches(triangle=False)!, \py!FedStatePattern!, 1.] A phaseless X state whose output is wired to input port $0$ of a W node with one input (or, with
\py!triangle=True!, of a $T$ or $T^{-1}$ with one input and one output); both joined ports claimed
exactly once; one shared dimension; both nodes outside every \bbox{}.
\py!FedStateMatch(state_id, target_id, wire, output_count, ...)!. Order (state, target).
Conditions for W: \code!state_is_a_phaseless_x_spider!, \code!target_is_a_w_node!,
\code!joined_and_unclaimed!, \code!same_dimension!, \code!outside_every_bang_box!; for triangles
\code!target_is_a_triangle! replaces the second.
\item[\py!find_triangle_effect_matches!, \py!TriangleEffectPattern!, 1.] A $T$ with one input and one output whose output feeds a phaseless X effect; joined ports claimed
exactly once; one shared dimension; outside every \bbox{}. Order (triangle, effect).
\item[\py!find_connective_inverse_matches(bind_first)!, \py!ConnectiveInversePattern!, 2.] Without \py!bind_first!: an S whose outputs $0$ and $1$ feed inputs $0$ and $1$ of a B. With
\py!bind_first!: a B whose output $0$ feeds input $0$ of an S. Each joining wire must join exactly
those ports, each claimed exactly once, with equal dimensions; the outer pairs (input $k$ of the
first with output $k$ of the second, for $k\in\{0,1\}$ when \py!bind_first!, $k=0$ otherwise)
have equal dimensions. Each outer pair yields a \py!Splice(wire, surviving_ref, far_ref)!: both
ports claimed at most once; the second port's wire is consumed when present, else the first's; the
far port is off the pair and is not the surviving port's neighbor. The far ports are distinct; every
spliced equation is solvable; both nodes outside every \bbox{}. Order: first id (one candidate per
first node).
\item[\py!find_connective_states_matches(swapped)!, \py!ConnectiveStatesPattern!, 1.] A B whose two inputs are each fed by a distinct phaseless Z state, or an S whose two outputs each
feed a distinct phaseless Z effect (X when swapped). Every connective port and every far port is
claimed exactly once by its wire; each end's shared dimension equals its connective port's
dimension; the connective's joint side has exactly one port; all three nodes outside every \bbox{}.
\py!ConnectiveStatesMatch(connective_id, end_ids, wires, joint_dim, effects, swapped, ...)!. Order:
connective id.
\item[\py!find_wz_bialgebra_matches!, \py!WZBialgebraPattern!, 1.] A phaseless Z spider with $n\ge 2$ inputs and one output whose output feeds the input of a W node
with one input and $m\ge 2$ outputs; joined ports claimed exactly once; one shared dimension;
outside every \bbox{}. \py!WZBialgebraMatch(z_id, w_id, wire, input_count=n, output_count=m, ...)!.
Order (Z, W).
\item[\py!find_w_effect_matches(zero)!, \py!WEffectPattern!, 1.] Output port $p$ of a W node with one input and at least two outputs (at least one when
\py!zero!) wired to a phaseless Z effect (an X effect when \py!zero!); joined ports claimed exactly
once; one shared dimension; outside every \bbox{}.
\py!WEffectMatch(w_id, effect_id, position=p, wire, output_count, zero, ...)!. Order (W, $p$).
\end{description}

The rules these patterns serve are: W fusion (two W nodes are one), W identity (a one-output W is
a bare wire), W zero copy (a phaseless X state copies through a W onto every output), triangle zero
state (a phaseless X state through $T$ or $T^{-1}$ is unchanged), triangle zero effect ($T$ into a
phaseless X effect is a phaseless Z effect), connective inverse (S then B, or B then S, is the
identity on the outer ports), connective states (two phaseless states into a B, or an S into two
effects, are one state or effect on the joint port), the W/Z bialgebra (a phaseless $n$-input Z
feeding an $m$-output W becomes $n$ W nodes and $m$ Z spiders), and the W effect (a W output into
a phaseless Z effect is a $T$ before the remaining W; into an X effect, the remaining W alone).
Their builders and scalars are specified in \cref{ch:rules}.

%======================================================================
\section{Port-scoped !-box matchers}\label{sec:rewriting-box}

The module \module{archytaszx.rewrite.match_box} holds the matcher for \term{port-box unfusion}:
a Z or X spider's port that alone forms the scope of a port-scoped \bbox{} moves onto a new
phaseless spider of the same color, joined to the rest by one wire.

\begin{signature}
\py!find_port_box_unfusion_matches(diagram, *, anchors=None) -> tuple[PortBoxUnfusionMatch, ...]!
\end{signature}
For every \bbox{} in ascending id order, a match requires: the box is port-scoped, has no parent,
and is no box's parent; its port scope is exactly one port $p$; $p$'s node is a registered Z or X
spider within one hop of an anchor; the spider has at least three ports; no other \bbox{} has $p$
in its port scope; the spider has a uniform port dimension; and the spider lies in no node-scoped
\bbox{}. \py!PortBoxUnfusionMatch(node_id, box_id, boxed=p, shared_dim, ...)! has support
\py!(node_id,)!. Matches are ordered by (node id, box id); \py!PortBoxUnfusionPattern! has radius
$1$. Its conditions are \code!node_is_a_spider!, \code!one_leg_port_boxed!,
\code!two_other_legs!, \code!same_dimension!, and \code!outside_every_bang_box!.

%======================================================================
\section{Matching across !-box boundaries}\label{sec:rewriting-boxcond}

A rewrite is sound on a diagram with \bboxes{} (\cref{ch:diagrams}) only when it commutes with
every instantiation of the \bboxes{}~\cite{dixon2009graphical,kissinger2014pattern}. Instead of
matching inside the instantiated diagram, every matcher checks a subset of five \term{box
conditions} on the symbolic diagram. \Cref{tab:rewriting-boxcond} states each condition as the code
enforces it; \cref{tab:rewriting-boxrules} lists which matcher enforces which.

\begin{table}[ht]
\centering\small
\caption{Box conditions, as the matchers enforce them.}\label{tab:rewriting-boxcond}
\begin{tabularx}{\textwidth}{@{}lX@{}}
\toprule
Condition & Requirement on a match \\
\midrule
O & No matched node lies in a node-scoped \bbox{}, has a port in a port-scoped \bbox{}, or is
wired to a node whose set of node boxes differs from its own (\emph{outside every \bbox{}}). \\
S & All matched nodes have the same innermost node-scoped \bbox{} (smallest node scope, ties by
id), or none lies in a node-scoped \bbox{}. \\
N & No node-scoped \bbox{} has a nonempty node scope contained in the set of nodes the
replacement deletes. \\
P & No port of the removed node lies in a port-scoped \bbox{}. \\
L & The splice keeps copies paired (\cref{sec:rewriting-identity}): no neighbor's boxes (minus the
link boxes of the far port) exceed the removed node's, the two far ports share no link box, and,
when the node is boxed, some wired neighbor has exactly the node's boxes; or the node's box
lifts to a port-scoped \bbox{} on the far port. \\
\bottomrule
\end{tabularx}
\end{table}

\begin{table}[ht]
\centering\small
\caption{Box conditions per matcher.}\label{tab:rewriting-boxrules}
\begin{tabularx}{\textwidth}{@{}>{\raggedright\arraybackslash}p{0.42\textwidth}>{\raggedright\arraybackslash}X@{}}
\toprule
Matcher & Conditions \\
\midrule
Spider fusion & S \\
$F^4$ cancellation & all four nodes have equal sets of node boxes \\
Cap, state copy, Hopf, bialgebra & O \\
Identity removal & P, L, and N unless the box lifts \\
Triangle--inverse cancellation & S, N, L (without link-port handling of own ports) \\
Fourier box on a state & S \\
W fusion & S, N \\
W identity & N, P, L (no lift) \\
Fed state, triangle effect, connective inverse and states, W/Z bialgebra, W effect & O \\
Port-box unfusion & spider in no node-scoped \bbox{}; the box top-level and childless \\
Pattern-diagram rules & no matched node in a node scope or with a port in a port scope, and no
matched node wired to such a node \\
\bottomrule
\end{tabularx}
\end{table}

The pattern-diagram condition is stronger than O: it excludes a wire to any boxed node, including
one boxed only through a port-scoped \bbox{}. The \bbox{} conditions of the axiom set and of the
discard rules are listed with each rule in \cref{ch:rules,ch:mixed}.


\begin{example}[Box conditions at work]\label{ex:rewriting-boxcond}
\Cref{fig:rewriting-boxcond} shows three situations. In (a), spider $u$ lies in a node-scoped
\bbox{} $b$ and spider $v$ outside it; their innermost boxes differ, so condition S fails and fusion
does not match: expanding $b$ at multiplicity $k$ would give $k$ copies of $u$, each wired to the
one $v$, and a single fusion cannot stand for $k$ of them. In (b), identity $u$ is alone in $b$, its
upper port is a boundary entry, and its lower port is wired to a spider $v$ outside $b$; the
identity is removed and $b$ \emph{lifts} to a port-scoped \bbox{} on the port of $v$ (condition L,
lift case), so that every instance has $k$ boundary ports on $v$. In (c), $u$ and $v$ share the
innermost box $b$; fusion matches (condition S), and each instance fuses $k$ disjoint pairs.
\end{example}

\begin{figure}[htbp]
\centering
\begin{minipage}[b]{0.3\textwidth}\centering
\begin{tikzpicture}[zxd]
  \node[zs] (u) at (0,1.6) {$u$}; \node[zs] (v) at (0,0.4) {$v$};
  \draw[wire] (v) -- (u); \draw[wire] (u) -- (0,2.5); \draw[wire] (0,-0.4) -- (v);
  \node[bang,fit=(u)] (b) {}; \node[banglabel] at (b.north east) {$b$};
\end{tikzpicture}\\(a) S fails: no fusion
\end{minipage}\hfill
\begin{minipage}[b]{0.3\textwidth}\centering
\begin{tikzpicture}[zxd]
  \node[zdot] (u) at (0,1.6) {}; \node[zs] (v) at (0,0.4) {$v$};
  \draw[wire] (v) -- (u); \draw[wire] (u) -- (0,2.5); \draw[wire] (-0.4,-0.4) -- (v);
  \draw[wire] (0.4,-0.4) -- (v);
  \node[bang,fit=(u)] (b) {}; \node[banglabel] at (b.north east) {$b$};
  \node at (0.9,1.0) {$\mapsto$};
  \node[zs] (w) at (1.8,0.4) {$v$};
  \draw[wire] (w) -- (1.8,2.5); \draw[wire] (1.4,-0.4) -- (w); \draw[wire] (2.2,-0.4) -- (w);
  \draw[bang] (1.6,1.1) rectangle (2.0,2.1); \node[banglabel] at (2.0,2.1) {$b$};
\end{tikzpicture}\\(b) L, lift case
\end{minipage}\hfill
\begin{minipage}[b]{0.3\textwidth}\centering
\begin{tikzpicture}[zxd]
  \node[zs] (u) at (0,1.6) {$u$}; \node[zs] (v) at (0,0.6) {$v$};
  \draw[wire] (v) -- (u); \draw[wire] (u) -- (0,2.6); \draw[wire] (0,-0.2) -- (v);
  \node[bang,fit=(u)(v)] (b) {}; \node[banglabel] at (b.north east) {$b$};
\end{tikzpicture}\\(c) S holds: fusion
\end{minipage}
\caption{Fusion and identity removal near a node-scoped \bbox{} $b$ (dashed). Wires that cross
the dashed outline end on a variable-arity side of a spider.}
\label{fig:rewriting-boxcond}
\end{figure}

\subsection{Declared side conditions}\label{sec:rewriting-declared}

Every fixed-shape pattern declares its side conditions as a module constant. The names and
descriptions below are the exact data a certificate carries; the order within each constant is
the order in which a match reports its outcomes. A swapped constant exchanges the colors of its
first entries and shares the rest. A template axiom reports \py!TEMPLATE_SIDE_CONDITIONS! followed
by its own conditions (\cref{sec:rewriting-template}).

\begin{xltabular}{\textwidth}{@{}>{\raggedright\arraybackslash}p{0.42\textwidth}>{\raggedright\arraybackslash}X@{}}
\caption{Declared side conditions of every fixed-shape and template pattern.}\label{tab:rewriting-declared}\\
\toprule
Constant and condition name & Description \\
\endfirsthead
\toprule
Constant and condition name & Description \\
\endhead
\multicolumn{2}{@{}l}{\texttt{BIALGEBRA\_SIDE\_CONDITIONS} (\file{match.py})}\\
\quad\texttt{x\_\allowbreak{}is\_\allowbreak{}a\_\allowbreak{}phaseless\_\allowbreak{}two\_\allowbreak{}to\_\allowbreak{}one} & an X spider with two inputs, one output, and no phase \\
\quad\texttt{z\_\allowbreak{}is\_\allowbreak{}a\_\allowbreak{}phaseless\_\allowbreak{}one\_\allowbreak{}to\_\allowbreak{}two} & a Z spider with one input, two outputs, and no phase \\
\quad\texttt{joined\_\allowbreak{}and\_\allowbreak{}unclaimed} & one wire joins them and neither port is otherwise used \\
\quad\texttt{same\_\allowbreak{}dimension} & every leg of the pair carries one dimension \\
\quad\texttt{outside\_\allowbreak{}every\_\allowbreak{}bang\_\allowbreak{}box} & neither node lies in a node-scope bang box, has a port-boxed leg, or is wired into one \\
\midrule
\multicolumn{2}{@{}l}{\texttt{BIALGEBRA\_SWAPPED\_SIDE\_CONDITIONS} (\file{match.py})}\\
\quad\texttt{z\_\allowbreak{}is\_\allowbreak{}a\_\allowbreak{}phaseless\_\allowbreak{}two\_\allowbreak{}to\_\allowbreak{}one} & a Z spider with two inputs, one output, and no phase \\
\quad\texttt{x\_\allowbreak{}is\_\allowbreak{}a\_\allowbreak{}phaseless\_\allowbreak{}one\_\allowbreak{}to\_\allowbreak{}two} & an X spider with one input, two outputs, and no phase \\
\quad\texttt{joined\_\allowbreak{}and\_\allowbreak{}unclaimed} & one wire joins them and neither port is otherwise used \\
\quad\texttt{same\_\allowbreak{}dimension} & every leg of the pair carries one dimension \\
\quad\texttt{outside\_\allowbreak{}every\_\allowbreak{}bang\_\allowbreak{}box} & neither node lies in a node-scope bang box, has a port-boxed leg, or is wired into one \\
\midrule
\multicolumn{2}{@{}l}{\texttt{CAP\_SIDE\_CONDITIONS} (\file{match.py})}\\
\quad\texttt{state\_\allowbreak{}is\_\allowbreak{}a\_\allowbreak{}z\_\allowbreak{}spider} & a Z spider with no input and one output \\
\quad\texttt{effect\_\allowbreak{}is\_\allowbreak{}an\_\allowbreak{}x\_\allowbreak{}spider} & an X spider with one input and no output \\
\quad\texttt{joined\_\allowbreak{}and\_\allowbreak{}unclaimed} & one wire joins them and neither port is otherwise used \\
\quad\texttt{phaseless} & neither node carries a phase vector \\
\quad\texttt{same\_\allowbreak{}dimension} & both legs carry one dimension \\
\quad\texttt{outside\_\allowbreak{}every\_\allowbreak{}bang\_\allowbreak{}box} & neither node lies in a node-scope bang box, has a port-boxed leg, or is wired into one \\
\midrule
\multicolumn{2}{@{}l}{\texttt{CAP\_SWAPPED\_SIDE\_CONDITIONS} (\file{match.py})}\\
\quad\texttt{state\_\allowbreak{}is\_\allowbreak{}an\_\allowbreak{}x\_\allowbreak{}spider} & an X spider with no input and one output \\
\quad\texttt{effect\_\allowbreak{}is\_\allowbreak{}a\_\allowbreak{}z\_\allowbreak{}spider} & a Z spider with one input and no output \\
\quad\texttt{joined\_\allowbreak{}and\_\allowbreak{}unclaimed} & one wire joins them and neither port is otherwise used \\
\quad\texttt{phaseless} & neither node carries a phase vector \\
\quad\texttt{same\_\allowbreak{}dimension} & both legs carry one dimension \\
\quad\texttt{outside\_\allowbreak{}every\_\allowbreak{}bang\_\allowbreak{}box} & neither node lies in a node-scope bang box, has a port-boxed leg, or is wired into one \\
\midrule
\multicolumn{2}{@{}l}{\texttt{FOURIER\_SIDE\_CONDITIONS} (\file{match.py})}\\
\quad\texttt{chain\_\allowbreak{}is\_\allowbreak{}series} & four F boxes joined output-to-input in series, each internal node used only there \\
\quad\texttt{same\_\allowbreak{}dimension} & every leg of the four boxes carries one dimension \\
\midrule
\multicolumn{2}{@{}l}{\texttt{FOURIER\_STATE\_SIDE\_CONDITIONS} (\file{match.py})}\\
\quad\texttt{spider\_\allowbreak{}is\_\allowbreak{}a\_\allowbreak{}phaseless\_\allowbreak{}z\_\allowbreak{}state} & a Z spider with no phase and a single leg \\
\quad\texttt{fourier\_\allowbreak{}box\_\allowbreak{}on\_\allowbreak{}that\_\allowbreak{}leg} & a registered F box wired to it in series \\
\quad\texttt{joined\_\allowbreak{}and\_\allowbreak{}unclaimed} & one wire joins them and neither port is otherwise used \\
\quad\texttt{free\_\allowbreak{}leg\_\allowbreak{}singly\_\allowbreak{}claimed} & the F box's other leg carries at most one claim \\
\quad\texttt{same\_\allowbreak{}dimension} & every leg of the pair carries one dimension \\
\quad\texttt{bang\_\allowbreak{}box\_\allowbreak{}scope\_\allowbreak{}agreement} & both nodes' innermost enclosing node-scope bang box, if any, are identical \\
\midrule
\multicolumn{2}{@{}l}{\texttt{FOURIER\_STATE\_SWAPPED\_SIDE\_CONDITIONS} (\file{match.py})}\\
\quad\texttt{spider\_\allowbreak{}is\_\allowbreak{}a\_\allowbreak{}phaseless\_\allowbreak{}x\_\allowbreak{}state} & an X spider with no phase and a single leg \\
\quad\texttt{fourier\_\allowbreak{}box\_\allowbreak{}on\_\allowbreak{}that\_\allowbreak{}leg} & a registered F box wired to it in series \\
\quad\texttt{joined\_\allowbreak{}and\_\allowbreak{}unclaimed} & one wire joins them and neither port is otherwise used \\
\quad\texttt{free\_\allowbreak{}leg\_\allowbreak{}singly\_\allowbreak{}claimed} & the F box's other leg carries at most one claim \\
\quad\texttt{same\_\allowbreak{}dimension} & every leg of the pair carries one dimension \\
\quad\texttt{bang\_\allowbreak{}box\_\allowbreak{}scope\_\allowbreak{}agreement} & both nodes' innermost enclosing node-scope bang box, if any, are identical \\
\midrule
\multicolumn{2}{@{}l}{\texttt{FUSION\_SIDE\_CONDITIONS} (\file{match.py})}\\
\quad\texttt{distinct\_\allowbreak{}nodes} & recorded fact, never a decision: the two matched nodes are not the same node \\
\quad\texttt{same\_\allowbreak{}generator\_\allowbreak{}type} & both nodes are the same registered spider color \\
\quad\texttt{parallel\_\allowbreak{}wires\_\allowbreak{}become\_\allowbreak{}self\_\allowbreak{}loops} & recorded fact, never a decision: every other wire joining the two nodes survives as a self-loop on the merged spider \\
\quad\texttt{consumed\_\allowbreak{}wire\_\allowbreak{}direction\_\allowbreak{}permitted\_\allowbreak{}for\_\allowbreak{}color} & for X, the consumed wire runs OUTPUT to INPUT; for Z, any direction combination is valid fusion \\
\quad\texttt{consumed\_\allowbreak{}ports\_\allowbreak{}singly\_\allowbreak{}claimed} & neither consumed port is claimed by a second wire, and neither is listed on either boundary list \\
\quad\texttt{bang\_\allowbreak{}box\_\allowbreak{}scope\_\allowbreak{}agreement} & both matched nodes' innermost enclosing node-scope bang box, if any, are identical (both unboxed, or the same box) -- Phase 7 \\
\quad\texttt{dimension\_\allowbreak{}agreement} & the connecting pair and every surviving leg of both nodes unify, in a bounded fixpoint, against the shared dimension -- equal outright, or unify defers or binds a symbol \\
\quad\texttt{phase\_\allowbreak{}dimension\_\allowbreak{}agreement} & every phase vector present unifies with the resolved shared leg dimension -- equal outright, or unify binds a symbol to a concrete value (never merely defers, and never binds to another still-symbolic Dim -- see the module docstring's 'Dimension constraints' note) \\
\quad\texttt{dimension\_\allowbreak{}guards\_\allowbreak{}satisfied} & every dimension guard the rule declares holds at the matched shared dimension \\
\midrule
\multicolumn{2}{@{}l}{\texttt{HOPF\_SIDE\_CONDITIONS} (\file{match.py})}\\
\quad\texttt{z\_\allowbreak{}and\_\allowbreak{}x\_\allowbreak{}spiders} & a registered Z spider and a registered X spider \\
\quad\texttt{two\_\allowbreak{}paths} & one Z output runs straight into an X input, a second Z output runs into that same X through two Fourier boxes in series \\
\quad\texttt{path\_\allowbreak{}ports\_\allowbreak{}unclaimed} & every port along the two paths carries exactly one wire \\
\quad\texttt{same\_\allowbreak{}dimension} & every leg along the two paths and both spiders share one dimension \\
\quad\texttt{phase\_\allowbreak{}dimension\_\allowbreak{}agreement} & each spider's phase vector, if present, is over that dimension \\
\quad\texttt{outside\_\allowbreak{}every\_\allowbreak{}bang\_\allowbreak{}box} & none of the four nodes lies in a node-scope bang box, has a port-boxed leg, or is wired into one \\
\midrule
\multicolumn{2}{@{}l}{\texttt{HOPF\_SWAPPED\_SIDE\_CONDITIONS} (\file{match.py})}\\
\quad\texttt{x\_\allowbreak{}and\_\allowbreak{}z\_\allowbreak{}spiders} & a registered X spider and a registered Z spider \\
\quad\texttt{two\_\allowbreak{}paths} & one X output runs straight into a Z input, a second X output runs into that same Z through two Fourier boxes in series \\
\quad\texttt{path\_\allowbreak{}ports\_\allowbreak{}unclaimed} & every port along the two paths carries exactly one wire \\
\quad\texttt{same\_\allowbreak{}dimension} & every leg along the two paths and both spiders share one dimension \\
\quad\texttt{phase\_\allowbreak{}dimension\_\allowbreak{}agreement} & each spider's phase vector, if present, is over that dimension \\
\quad\texttt{outside\_\allowbreak{}every\_\allowbreak{}bang\_\allowbreak{}box} & none of the four nodes lies in a node-scope bang box, has a port-boxed leg, or is wired into one \\
\midrule
\multicolumn{2}{@{}l}{\texttt{IDENTITY\_SIDE\_CONDITIONS} (\file{match.py})}\\
\quad\texttt{node\_\allowbreak{}is\_\allowbreak{}a\_\allowbreak{}phaseless\_\allowbreak{}spider} & a registered Z or X spider with no phase \\
\quad\texttt{one\_\allowbreak{}in\_\allowbreak{}one\_\allowbreak{}out} & exactly one input leg and one output leg \\
\quad\texttt{not\_\allowbreak{}a\_\allowbreak{}self\_\allowbreak{}loop} & no wire joins the node's two legs \\
\quad\texttt{legs\_\allowbreak{}singly\_\allowbreak{}claimed} & neither leg carries a second wire or a boundary slot alongside one, and at least one leg carries a wire to splice through \\
\quad\texttt{splice\_\allowbreak{}keeps\_\allowbreak{}two\_\allowbreak{}ports} & the spliced wire's far port is not the surviving leg's own neighbor \\
\quad\texttt{same\_\allowbreak{}dimension} & both legs carry one dimension \\
\quad\texttt{leaves\_\allowbreak{}every\_\allowbreak{}bang\_\allowbreak{}box\_\allowbreak{}populated} & no node-scope bang box holds the node and nothing else, unless that box lifts to a port-scope box \\
\quad\texttt{no\_\allowbreak{}port\_\allowbreak{}boxed\_\allowbreak{}leg} & neither leg lies in a port-scope bang box \\
\quad\texttt{splice\_\allowbreak{}keeps\_\allowbreak{}copies} & no neighbor lies in a node-scope bang box nested inside the node's, and some wired neighbor shares the node's boxes, unless the node's box lifts to a port-scope box: it holds only the node, has no child box, the surviving leg is a boundary slot, and the far port is a spider leg one box out with no port-boxed leg \\
\midrule
\multicolumn{2}{@{}l}{\texttt{STATE\_COPY\_SIDE\_CONDITIONS} (\file{match.py})}\\
\quad\texttt{state\_\allowbreak{}is\_\allowbreak{}a\_\allowbreak{}phaseless\_\allowbreak{}x\_\allowbreak{}spider} & an X spider with no input, one output, and no phase \\
\quad\texttt{spider\_\allowbreak{}is\_\allowbreak{}a\_\allowbreak{}phaseless\_\allowbreak{}z\_\allowbreak{}spider} & a Z spider with one input and no phase \\
\quad\texttt{joined\_\allowbreak{}and\_\allowbreak{}unclaimed} & one wire joins them and neither port is otherwise used \\
\quad\texttt{same\_\allowbreak{}dimension} & every leg of the pair carries one dimension \\
\quad\texttt{outside\_\allowbreak{}every\_\allowbreak{}bang\_\allowbreak{}box} & neither node lies in a node-scope bang box, has a port-boxed leg, or is wired into one \\
\midrule
\multicolumn{2}{@{}l}{\texttt{STATE\_COPY\_SWAPPED\_SIDE\_CONDITIONS} (\file{match.py})}\\
\quad\texttt{state\_\allowbreak{}is\_\allowbreak{}a\_\allowbreak{}phaseless\_\allowbreak{}z\_\allowbreak{}spider} & a Z spider with no input, one output, and no phase \\
\quad\texttt{spider\_\allowbreak{}is\_\allowbreak{}a\_\allowbreak{}phaseless\_\allowbreak{}x\_\allowbreak{}spider} & an X spider with one input and no phase \\
\quad\texttt{joined\_\allowbreak{}and\_\allowbreak{}unclaimed} & one wire joins them and neither port is otherwise used \\
\quad\texttt{same\_\allowbreak{}dimension} & every leg of the pair carries one dimension \\
\quad\texttt{outside\_\allowbreak{}every\_\allowbreak{}bang\_\allowbreak{}box} & neither node lies in a node-scope bang box, has a port-boxed leg, or is wired into one \\
\midrule
\multicolumn{2}{@{}l}{\texttt{TRIANGLE\_SIDE\_CONDITIONS} (\file{match.py})}\\
\quad\texttt{inverse\_\allowbreak{}pair\_\allowbreak{}in\_\allowbreak{}series} & a T and a Ti joined output-to-input, either order \\
\quad\texttt{joining\_\allowbreak{}wire\_\allowbreak{}unclaimed} & neither joined port carries a second claim \\
\quad\texttt{outer\_\allowbreak{}legs\_\allowbreak{}singly\_\allowbreak{}claimed} & neither outer leg carries a second claim, and at least one carries a wire to splice through \\
\quad\texttt{splice\_\allowbreak{}keeps\_\allowbreak{}two\_\allowbreak{}ports} & the spliced wire's far port is not the surviving outer leg's own neighbor \\
\quad\texttt{same\_\allowbreak{}dimension} & every leg of the pair carries one dimension \\
\quad\texttt{bang\_\allowbreak{}box\_\allowbreak{}scope\_\allowbreak{}agreement} & both nodes' innermost enclosing node-scope bang box, if any, are identical \\
\quad\texttt{leaves\_\allowbreak{}every\_\allowbreak{}bang\_\allowbreak{}box\_\allowbreak{}populated} & no node-scope bang box holds the pair and nothing else \\
\midrule
\multicolumn{2}{@{}l}{\texttt{CONNECTIVE\_INVERSE\_SIDE\_CONDITIONS} (\file{match_zxw.py})}\\
\quad\texttt{inverse\_\allowbreak{}pair} & an S whose two outputs feed a B's two inputs in order, or a B whose output feeds an S \\
\quad\texttt{joining\_\allowbreak{}wires\_\allowbreak{}unclaimed} & no joined port carries a second claim \\
\quad\texttt{matching\_\allowbreak{}dimensions} & each joined pair of legs, and each outer pair, carry one dimension \\
\quad\texttt{outer\_\allowbreak{}legs\_\allowbreak{}singly\_\allowbreak{}claimed} & no outer leg carries a second claim, and each identity to splice has a wire \\
\quad\texttt{splice\_\allowbreak{}keeps\_\allowbreak{}two\_\allowbreak{}ports} & no spliced wire's far port is on the pair or is the surviving leg's own neighbor \\
\quad\texttt{outside\_\allowbreak{}every\_\allowbreak{}bang\_\allowbreak{}box} & neither node lies in a node-scope bang box, has a port-boxed leg, or is wired into one \\
\midrule
\multicolumn{2}{@{}l}{\texttt{CONNECTIVE\_STATES\_SIDE\_CONDITIONS} (\file{match_zxw.py})}\\
\quad\texttt{connective\_\allowbreak{}with\_\allowbreak{}two\_\allowbreak{}ends} & a B whose two inputs are fed by phaseless Z states, or an S whose two outputs feed phaseless Z effects \\
\quad\texttt{joined\_\allowbreak{}and\_\allowbreak{}unclaimed} & each joining wire's ports carry no second claim \\
\quad\texttt{matching\_\allowbreak{}dimensions} & each state or effect carries its connective leg's dimension \\
\quad\texttt{outside\_\allowbreak{}every\_\allowbreak{}bang\_\allowbreak{}box} & no node lies in a node-scope bang box, has a port-boxed leg, or is wired into one \\
\midrule
\multicolumn{2}{@{}l}{\texttt{CONNECTIVE\_STATES\_SWAPPED\_SIDE\_CONDITIONS} (\file{match_zxw.py})}\\
\quad\texttt{connective\_\allowbreak{}with\_\allowbreak{}two\_\allowbreak{}ends} & a B whose two inputs are fed by phaseless X states, or an S whose two outputs feed phaseless X effects \\
\quad\texttt{joined\_\allowbreak{}and\_\allowbreak{}unclaimed} & each joining wire's ports carry no second claim \\
\quad\texttt{matching\_\allowbreak{}dimensions} & each state or effect carries its connective leg's dimension \\
\quad\texttt{outside\_\allowbreak{}every\_\allowbreak{}bang\_\allowbreak{}box} & no node lies in a node-scope bang box, has a port-boxed leg, or is wired into one \\
\midrule
\multicolumn{2}{@{}l}{\texttt{TRIANGLE\_EFFECT\_SIDE\_CONDITIONS} (\file{match_zxw.py})}\\
\quad\texttt{node\_\allowbreak{}is\_\allowbreak{}a\_\allowbreak{}triangle} & a registered T with one input and one output \\
\quad\texttt{effect\_\allowbreak{}is\_\allowbreak{}a\_\allowbreak{}phaseless\_\allowbreak{}x\_\allowbreak{}spider} & an X spider with one input, no output, and no phase \\
\quad\texttt{joined\_\allowbreak{}and\_\allowbreak{}unclaimed} & one wire joins them and neither port is otherwise used \\
\quad\texttt{same\_\allowbreak{}dimension} & every leg of the pair carries one dimension \\
\quad\texttt{outside\_\allowbreak{}every\_\allowbreak{}bang\_\allowbreak{}box} & neither node lies in a node-scope bang box, has a port-boxed leg, or is wired into one \\
\midrule
\multicolumn{2}{@{}l}{\texttt{TRIANGLE\_STATE\_SIDE\_CONDITIONS} (\file{match_zxw.py})}\\
\quad\texttt{state\_\allowbreak{}is\_\allowbreak{}a\_\allowbreak{}phaseless\_\allowbreak{}x\_\allowbreak{}spider} & an X spider with no input, one output, and no phase \\
\quad\texttt{target\_\allowbreak{}is\_\allowbreak{}a\_\allowbreak{}triangle} & a registered T or Ti, fed at its input \\
\quad\texttt{joined\_\allowbreak{}and\_\allowbreak{}unclaimed} & one wire joins them and neither port is otherwise used \\
\quad\texttt{same\_\allowbreak{}dimension} & every leg of the pair carries one dimension \\
\quad\texttt{outside\_\allowbreak{}every\_\allowbreak{}bang\_\allowbreak{}box} & neither node lies in a node-scope bang box, has a port-boxed leg, or is wired into one \\
\midrule
\multicolumn{2}{@{}l}{\texttt{WZ\_BIALGEBRA\_SIDE\_CONDITIONS} (\file{match_zxw.py})}\\
\quad\texttt{z\_\allowbreak{}is\_\allowbreak{}a\_\allowbreak{}phaseless\_\allowbreak{}fan\_\allowbreak{}in} & a phaseless Z spider with at least two inputs and exactly one output \\
\quad\texttt{w\_\allowbreak{}is\_\allowbreak{}a\_\allowbreak{}fan\_\allowbreak{}out} & a registered W node with at least two outputs \\
\quad\texttt{joining\_\allowbreak{}wire\_\allowbreak{}unclaimed} & neither joined port carries a second claim \\
\quad\texttt{same\_\allowbreak{}dimension} & every leg of both nodes carries one dimension \\
\quad\texttt{outside\_\allowbreak{}every\_\allowbreak{}bang\_\allowbreak{}box} & neither node lies in a node-scope bang box, has a port-boxed leg, or is wired into one \\
\midrule
\multicolumn{2}{@{}l}{\texttt{W\_COPY\_SIDE\_CONDITIONS} (\file{match_zxw.py})}\\
\quad\texttt{state\_\allowbreak{}is\_\allowbreak{}a\_\allowbreak{}phaseless\_\allowbreak{}x\_\allowbreak{}spider} & an X spider with no input, one output, and no phase \\
\quad\texttt{target\_\allowbreak{}is\_\allowbreak{}a\_\allowbreak{}w\_\allowbreak{}node} & a registered W node, fed at its input \\
\quad\texttt{joined\_\allowbreak{}and\_\allowbreak{}unclaimed} & one wire joins them and neither port is otherwise used \\
\quad\texttt{same\_\allowbreak{}dimension} & every leg of the pair carries one dimension \\
\quad\texttt{outside\_\allowbreak{}every\_\allowbreak{}bang\_\allowbreak{}box} & neither node lies in a node-scope bang box, has a port-boxed leg, or is wired into one \\
\midrule
\multicolumn{2}{@{}l}{\texttt{W\_FUSION\_SIDE\_CONDITIONS} (\file{match_zxw.py})}\\
\quad\texttt{two\_\allowbreak{}w\_\allowbreak{}nodes} & a W output wired into a second W's input \\
\quad\texttt{joining\_\allowbreak{}wire\_\allowbreak{}unclaimed} & neither joined port carries a second claim \\
\quad\texttt{same\_\allowbreak{}dimension} & every leg of both nodes carries one dimension \\
\quad\texttt{bang\_\allowbreak{}box\_\allowbreak{}scope\_\allowbreak{}agreement} & both nodes' innermost enclosing node-scope bang box, if any, are identical \\
\quad\texttt{leaves\_\allowbreak{}every\_\allowbreak{}bang\_\allowbreak{}box\_\allowbreak{}populated} & no node-scope bang box holds the pair and nothing else \\
\midrule
\multicolumn{2}{@{}l}{\texttt{W\_IDENTITY\_SIDE\_CONDITIONS} (\file{match_zxw.py})}\\
\quad\texttt{node\_\allowbreak{}is\_\allowbreak{}a\_\allowbreak{}w\_\allowbreak{}node} & a registered W node \\
\quad\texttt{one\_\allowbreak{}in\_\allowbreak{}one\_\allowbreak{}out} & exactly one input leg and one output leg \\
\quad\texttt{not\_\allowbreak{}a\_\allowbreak{}self\_\allowbreak{}loop} & no wire joins the node's two legs \\
\quad\texttt{legs\_\allowbreak{}singly\_\allowbreak{}claimed} & neither leg carries a second wire or a boundary slot alongside one, and at least one leg carries a wire to splice through \\
\quad\texttt{splice\_\allowbreak{}keeps\_\allowbreak{}two\_\allowbreak{}ports} & the spliced wire's far port is not the surviving leg's own neighbor \\
\quad\texttt{same\_\allowbreak{}dimension} & both legs carry one dimension \\
\quad\texttt{leaves\_\allowbreak{}every\_\allowbreak{}bang\_\allowbreak{}box\_\allowbreak{}populated} & no node-scope bang box holds the node and nothing else \\
\quad\texttt{no\_\allowbreak{}port\_\allowbreak{}boxed\_\allowbreak{}leg} & neither leg lies in a port-scope bang box \\
\midrule
\multicolumn{2}{@{}l}{\texttt{W\_ZERO\_EFFECT\_SIDE\_CONDITIONS} (\file{match_zxw.py})}\\
\quad\texttt{node\_\allowbreak{}is\_\allowbreak{}a\_\allowbreak{}w\_\allowbreak{}node} & a registered W node with at least one output \\
\quad\texttt{effect\_\allowbreak{}is\_\allowbreak{}a\_\allowbreak{}phaseless\_\allowbreak{}x\_\allowbreak{}spider} & an X spider with one input, no output, and no phase \\
\quad\texttt{joined\_\allowbreak{}and\_\allowbreak{}unclaimed} & one wire joins them and neither port is otherwise used \\
\quad\texttt{same\_\allowbreak{}dimension} & every leg of the pair carries one dimension \\
\quad\texttt{outside\_\allowbreak{}every\_\allowbreak{}bang\_\allowbreak{}box} & neither node lies in a node-scope bang box, has a port-boxed leg, or is wired into one \\
\midrule
\multicolumn{2}{@{}l}{\texttt{W\_Z\_EFFECT\_SIDE\_CONDITIONS} (\file{match_zxw.py})}\\
\quad\texttt{node\_\allowbreak{}is\_\allowbreak{}a\_\allowbreak{}w\_\allowbreak{}node} & a registered W node with at least two outputs \\
\quad\texttt{effect\_\allowbreak{}is\_\allowbreak{}a\_\allowbreak{}phaseless\_\allowbreak{}z\_\allowbreak{}spider} & a Z spider with one input, no output, and no phase \\
\quad\texttt{joined\_\allowbreak{}and\_\allowbreak{}unclaimed} & one wire joins them and neither port is otherwise used \\
\quad\texttt{same\_\allowbreak{}dimension} & every leg of the pair carries one dimension \\
\quad\texttt{outside\_\allowbreak{}every\_\allowbreak{}bang\_\allowbreak{}box} & neither node lies in a node-scope bang box, has a port-boxed leg, or is wired into one \\
\midrule
\multicolumn{2}{@{}l}{\texttt{PORT\_BOX\_UNFUSION\_SIDE\_CONDITIONS} (\file{match_box.py})}\\
\quad\texttt{node\_\allowbreak{}is\_\allowbreak{}a\_\allowbreak{}spider} & a registered Z or X spider \\
\quad\texttt{one\_\allowbreak{}leg\_\allowbreak{}port\_\allowbreak{}boxed} & a top-level port-scope bang box with no children holds exactly this one leg \\
\quad\texttt{two\_\allowbreak{}other\_\allowbreak{}legs} & the spider has at least two legs besides the boxed one \\
\quad\texttt{same\_\allowbreak{}dimension} & every leg of the spider carries one dimension \\
\quad\texttt{outside\_\allowbreak{}every\_\allowbreak{}bang\_\allowbreak{}box} & the spider lies in no node-scope bang box \\
\midrule
\multicolumn{2}{@{}l}{\texttt{TEMPLATE\_SIDE\_CONDITIONS} (\file{match_template.py})}\\
\quad\texttt{template\_\allowbreak{}matches} & the side's template matches up to leg symmetry \\
\quad\texttt{parameters\_\allowbreak{}match} & every node's parameter matches its template \\
\quad\texttt{same\_\allowbreak{}dimension} & every matched leg carries the anchor's dimension \\
\quad\texttt{outside\_\allowbreak{}every\_\allowbreak{}bang\_\allowbreak{}box} & no matched node or leg lies in a bang box \\
\quad\texttt{right\_\allowbreak{}side\_\allowbreak{}defined} & every parameter of the built side is defined \\
\bottomrule
\end{xltabular}

%======================================================================
\section{The rewrite engine}\label{sec:rewriting-engine}

The module \module{archytaszx.rewrite.engine} applies one rule at one match and drives repeated
application. \py!apply! is generic over every rule: it reads only the \py!Match! protocol and the
\py!BuildResult! record, never a match's rule-specific fields. The module evaluates nothing
numerically.

\subsection{Applying a rule}

\begin{signature}
\py!apply(diagram: Diagram, rule: Rule, match: Match, *, input_report: ValidationReport | None = None) -> RewriteResult!
\end{signature}
\py!apply! never mutates \py!diagram!. \py!input_report!, when given, is taken as
\py!validate(diagram)! (\cref{ch:diagrams}). It returns a frozen
\py!RewriteResult(diagram, new_node_ids, step, report)!, where \py!report! is the validation
report of the new diagram (excluded from equality).

\begin{algorithm}[Rule application]\label{alg:rewriting-apply}
\leavevmode
\begin{enumerate}
\item \emph{Coverage.} Run \cref{alg:rewriting-coverage} with \py!rule.side_conditions! and context
  \py!rule.name!.
\item \emph{Copy.} Let \py!working = diagram.copy()!; the parameter environment rides along.
\item \emph{Build.} Call \py!rule.builder(working, match)!. Then update chain-\bbox{} links
  (below).
\item \emph{Check the build.} Raise \py!RewriteGrammarError! when \py!build_result.diagram! is not
  \py!working! by identity; when the wire set of \py!working! differs from that of \py!diagram!
  (the message lists the added and removed wires); when either boundary list differs. Raise
  \py!RewriteDomainError! when the builder's scalar differs from \py!rule.scalar_for_match(match)!.
  Raise \py!RewriteGrammarError!, in this order, when: a consumed wire or node is absent; a consumed
  wire is listed twice; a consumed node is listed twice; a new node id names no node, or a
  port-mapping value names no real port; a port-mapping value lies on a consumed node; the phase
  substitutions name an unconsumed node; a node that is neither consumed nor new differs from its
  original (an in-place edit); a new node id repeats; a new node id already existed in the input; a
  port-mapping key lies on no consumed node; the port mapping is not injective.
\item \emph{Remap.} For every wire of \py!working! in sorted order: drop it if consumed; otherwise,
  if an endpoint lies on a consumed node, replace the endpoint through the port mapping (an
  unmapped endpoint on a consumed node raises \py!RewriteDomainError!), raise
  \py!RewriteGrammarError! if both endpoints collapse to one port, and re-add the wire. Then add
  every new wire, raising \py!RewriteGrammarError! when an endpoint names no real port or lies on a
  consumed node. Require the wire count to equal (input wire count) $-$ (distinct consumed wires)
  $+$ (distinct new wires), else raise \py!RewriteGrammarError!. Rebuild both boundary lists
  through the same endpoint remap, position by position.
\item \emph{Remove.} Remove the consumed nodes in order; then check far link ends (below).
\item \emph{Scalar.} Multiply the scalar of \py!working! by the builder's scalar.
\item \emph{Relative validation.} Validate \py!working!. Translate every hard issue of the input
  into post-rewrite coordinates and count both sides as multisets of keys
  (\emph{kind}, \emph{reference}). If the result has a key with a count above the input's, raise
  \py!RewriteDomainError! (\code!rule '<name>': rewrite introduced hard-error issue kind(s) not present in the input diagram: <kind> at <ref> (xN); ...!), the entries sorted by (kind value,
  \py!repr! of reference). A diagram that already carries a hard error therefore stays rewritable.
  Deferred issues never block; the same multiset difference over them yields the removed and
  introduced deferred issues of the step.
\item \emph{Provenance.} Return the result with a \py!RewriteStep!.
\end{enumerate}
\end{algorithm}


\paragraph{Messages.}
Every message of \py!apply! begins with \code!rule '<name>': !. The checks of steps 4 and 5 produce
the following texts (abridged after the first clause, which is fixed):
\begin{center}\small
\begin{tabularx}{\textwidth}{@{}>{\raggedright\arraybackslash}p{0.28\textwidth}X@{}}
\toprule
Condition & Message after the rule prefix \\
\midrule
result not the working copy & \code!builder returned a BuildResult.diagram that is not the working diagram it was given; ...! \\
wire set edited & \code!builder mutated the working diagram's wires directly (added [...]; removed [...]); ...! \\
boundary edited & \code!builder mutated the working diagram's boundary lists directly (inputs .. -> ..; outputs .. -> ..); ...! \\
scalar mismatch & \code!declares scalar_introduced=<s> at this match's dimension, but its builder returned <s'>! \\
absent wire or node & \code!match does not belong to the diagram it is applied to (consumed wire(s) absent: [...]; consumed node id(s) absent: [...])! \\
repeated consumed wire & \code!build_result.consumed_wires names the same wire more than once: [...] -- ...! \\
repeated consumed node & \code!build_result.consumed_node_ids names the same node id more than once: [...] -- ...! \\
unreal new node or port & \code!builder-reported BuildResult fields do not name real nodes/ports in the working diagram (...)! \\
mapping onto a consumed node & \code!builder's port_mapping sends a surviving port onto (...), which lies on a node this rewrite consumes; ...! \\
foreign phase substitution & \code!build_result.phase_substitutions names node id(s) [...] that this rewrite does not consume; ...! \\
in-place edit & \code!builder edited the generator type, port dims or phase of pre-existing node(s) [...] in place; ...! \\
repeated new node & \code!build_result.new_node_ids names the same node id more than once: [...] -- ...! \\
new node pre-existing & \code!build_result.new_node_ids names node id(s) [...] that already existed in the input diagram; ...! \\
key off the match & \code!builder's port_mapping keys [...] lie on no consumed node, ...! \\
not injective & \code!builder's port_mapping is not injective -- ...! \\
unmapped endpoint & \code!port <ref> is on a consumed node but is absent from the builder's port_mapping; ...! \\
collapsed wire & \code!port_mapping collapses wire <w> onto a single port <p>; ...! \\
bad new wire & \code!build_result.new_wires names port <p>, which is not a real port in the working diagram! (or \code!... on a node this rewrite consumes!) \\
wire count & \code!wire-count postcondition violated after remapping -- expected N wire(s) (M before this rewrite minus K distinct consumed), got N'; ...! \\
\bottomrule
\end{tabularx}
\end{center}

\begin{example}[Remapping]\label{ex:rewriting-remap}
In \cref{ex:rewriting-fusion}, the build result of fusion consumes nodes $0$ and $1$ and the wire
between them, and maps $(0,\mathrm{in},0)\mapsto(2,\mathrm{in},0)$ and
$(1,\mathrm{out},0)\mapsto(2,\mathrm{out},0)$. Step 5 drops the consumed wire; no other wire
touches a consumed node; the wire count goes from $1$ to $1-1+0=0$; the boundary lists become
$((2,\mathrm{in},0))$ and $((2,\mathrm{out},0))$. Step 6 removes nodes $0$ and $1$.
\end{example}

\paragraph{Issue keys.}
The key of a validation issue is (kind, reference), where the reference is the issue's port, else
its wire, else its node id, else \py!None!; messages never enter the key. To translate an input
issue, a port on a consumed node goes through the port mapping (an unmapped one stays as is and
matches nothing); a wire translates both endpoints and falls back to the untranslated wire when
they collapse; a node id on a consumed node becomes the single new node id when exactly one node
was created, and otherwise stays unchanged. For the deferred issues, the surplus of each key is
taken from the issues in validation order; a step is \py!deferred_issue_identity_ambiguous! when
some key's surplus is positive but smaller than its number of issues.

\paragraph{Chain \bboxes{} and links.}
After the builder runs, \py!_follow_links! examines every \bbox{} of the input in id order, with
$G$ the consumed nodes. It raises \py!RewriteDomainError! when an extended box (\cref{ch:diagrams})
no longer exists in \py!working!; when a pair, tree, or fed box has a node in $G$; when the far end
of a tree box's parent port lies in $G$; when a left or right port of a pair box lies on $G$; when
the box has a node in $G$, carries an index symbol, and that symbol is free in the introduced
scalar; and when the far end of a link port lies in $G$ outside the box's node scope, unless the
rule is \code!identity_removal! or \code!spider_fusion! and the box has no node in $G$. For a
surviving box whose links touch $G$, each link port on $G$ is replaced by its image under the port
mapping; an image on $G$ or outside the box's node scope in \py!working! raises
\py!RewriteDomainError!. After node removal, a link port whose far end was consumed must be a
boundary slot or wired to a node off its box; otherwise \py!RewriteDomainError! is raised.

\begin{definition}[Rewrite step]\label{def:rewriting-step}
A \term{rewrite step} is the frozen provenance record
\begin{signature}
\py!RewriteStep(rule_name, match, consumed_node_ids, consumed_wires, side_condition_outcomes, dimension_constraints, scalar_introduced, port_mapping, new_node_ids, removed_deferred_issues=(), introduced_deferred_issues=(), phase_substitutions={}, deferred_issue_identity_ambiguous=False)!
\end{signature}
\py!match! is stored verbatim, so the step can be replayed by looking up \py!rule_name! and
applying at the same match without rematching. \py!side_condition_outcomes! and
\py!dimension_constraints! come from the match; \py!scalar_introduced!, \py!port_mapping!, and the
node and wire fields from the build result; the deferred-issue fields from step~8.
\py!phase_substitutions! is empty when the builder supplied \py!None!. The hash folds the port
mapping and the phase substitutions as unordered sets, so equal steps hash equally within one
process.
\end{definition}

A \code!DEFERRED! dimension constraint in a step records an assumption, not a satisfiable
equation: a surviving port of dimension $d^2$ forced onto a shared $d$ records $d^2 = d$, which
holds only at $d = 1$. The certificate layer (\cref{ch:proofs}) carries these constraints.

\subsection{Strategies and the termination guard}

\begin{definition}[Termination guard]\label{def:rewriting-guardrec}
\py!TerminationGuard(max_steps: int, max_nodes: int | None = None, detect_loops: bool = True)! is
the budget of a strategy loop. \py!max_steps! and \py!max_nodes! are non-negative integers that are
not \py!bool!; \py!detect_loops! is a \py!bool!; violations raise \py!RewriteGrammarError!.
\py!DEFAULT_GUARD = TerminationGuard(max_steps=128, max_nodes=None, detect_loops=True)!.
\end{definition}

\py!StopReason! has four members: \code!FIXPOINT! (no rule matched), \code!STEP_LIMIT!
(\py!max_steps! rewrites applied), \code!NODE_LIMIT! (more than \py!max_nodes! nodes), and
\code!LOOP_DETECTED! (a rewrite reproduced a visited diagram). A run returns a frozen
\py!StrategyOutcome(diagram, steps, stop_reason, scalar_accumulated, steps_attempted)!, where
\py!scalar_accumulated! is the exact product of the steps' scalars and \py!steps_attempted! counts
the searches for an applicable rule. \py!TerminationGuardTripped(outcome)!, a \py!RewriteError!,
carries the outcome a non-strict run would have returned; its message is
\code!strategy stopped at '<reason>' rather than a fixpoint after N step(s)!.

\begin{signature}
\py!apply_until_fixpoint(diagram, rules: Sequence[Rule], *, guard=DEFAULT_GUARD, strict=False, cache: RewriteCache | None = None, on_result=None) -> StrategyOutcome!
\end{signature}
raises \py!RewriteGrammarError! when \py!guard! is not a \py!TerminationGuard!, \py!cache! is
neither \py!None! nor a \py!RewriteCache!, or \py!on_result! is neither \py!None! nor callable.
\begin{algorithm}[Fixpoint strategy]\label{alg:rewriting-fixpoint}
Let the \emph{loop measure} of a diagram be the triple (wire count, node count, total output count
of W nodes). Initialize the history with the input's canonical key (\cref{ch:diagrams}) when loop
detection is on.
\begin{enumerate}
\item Count an attempt. If \py!max_nodes! is set and the current node count exceeds it, stop with
  \code!NODE_LIMIT!. If the step count has reached \py!max_steps!, stop with \code!STEP_LIMIT!.
\item Find the first applicable rule: for each rule in order, obtain its matches (below). When the
  diagram has no extended \bbox{}, the first match of the first rule with matches is taken.
  Otherwise each match is trial-applied in order and the first that does not raise
  \py!RewriteDomainError! is taken. None found: stop with \code!FIXPOINT!.
\item Apply it, passing the previous result's validation report as the input report; hand the
  result to \py!on_result!; append the step; multiply the scalar.
\item If loop detection is off or the new measure is lexicographically smaller than the previous
  one, continue. Otherwise look the new canonical key up in the history; if an isomorphic diagram is
  stored under it, stop with \code!LOOP_DETECTED!; else store the diagram under its key.
\end{enumerate}
With \py!strict=True!, any stop other than \code!FIXPOINT! raises \py!TerminationGuardTripped!.
\end{algorithm}
The canonical key reads the global scalar, so a cycle that multiplies a factor on each round does
not register as a loop and runs to the step limit.

\paragraph{Match sources.}
Without a cache, a rule's matches are \py!rule.pattern.find_matches(diagram)!. With a cache, each
diagram is fingerprinted once per iteration (\cref{sec:rewriting-cache}); when the cache carries an
incremental matcher that tracks the rule's pattern (the same object, or an equal pattern of the same
type with the same pattern key), the matches come from the matcher's rematch of this diagram;
otherwise from the cache's match memo. A \py!CacheError! from any cache operation is re-raised as
\py!RewriteGrammarError! with a prefix naming the operation.

\begin{signature}
\py!toward_normal_form(diagram, *, guard=DEFAULT_GUARD, strict=False, cache=None) -> StrategyOutcome!
\end{signature}
runs the fixpoint strategy over \py!normal_form_rules()!: the rules of \py!NORMAL_FORM_RULE_NAMES!
that \py!lookup_rule! (\cref{ch:rules}) resolves, in this order:
\code!identity_removal!, \code!triangle_inverse_cancellation!, \code!w_identity!, \code!w_fusion!,
\code!connective_inverse!, \code!connective_inverse_bind_first!, \code!connective_states!,
\code!connective_states_swapped!, \code!triangle_zero_state!, \code!triangle_zero_effect!,
\code!w_zero_effect!, \code!w_z_effect!, \code!fourier_cancellation!,
\code!fourier_state_color_change!, \code!fourier_state_color_change_swapped!, \code!spider_fusion!,
\code!hopf!, \code!hopf_swapped!, \code!state_copy!, \code!state_copy_swapped!, \code!w_zero_copy!,
\code!zx_cap!, \code!zx_cap_swapped!, \code!zxw_ept!, \code!scalar_box!.
\py!missing_normal_form_rule_names()! returns the names that do not resolve. Every rule of the list
except \code!scalar_box! and \code!w_z_effect! removes at least one wire; \code!scalar_box! removes a
node and no wire; \code!w_z_effect! keeps the wire and node counts and removes one W output. The
bialgebra rules are absent, since they turn two nodes into four.

%======================================================================
\section{Normal forms}\label{sec:rewriting-nf}

The module \module{archytaszx.rewrite.normal_form} reduces a diagram to the fixpoint of a rule set
and keys the result so that two reductions can be compared. A normal form here is the output of a
terminating strategy, not a unique representative: the procedure is sound, and two diagrams with
the same key are equal, but two equal diagrams can reach different keys.

\begin{definition}[Normal form record]\label{def:rewriting-nfrec}
\py!NormalForm(source: Diagram, diagram: Diagram, outcome: StrategyOutcome, results: tuple[RewriteResult, ...], key: str)!
is a frozen record compared by identity. \py!source! is a copy of the input, \py!diagram! the exact
strategy output (so the recorded results replay onto it id for id), \py!results! every rewrite
result in application order, and \py!key! the view key of the comparison view of \py!diagram!.
Construction raises \py!RewriteGrammarError! for a field of the wrong type, checked in declaration
order. \py!reached_fixpoint! is true iff the stop reason is \code!FIXPOINT!;
\py!assumed_constraints! is \py!assumed! (\cref{sec:rewriting-rule}) over the dimension
constraints of every step.
\end{definition}

\begin{signature}
\py!normal_form(diagram, *, rules: Sequence[Rule] | None = None, guard=DEFAULT_GUARD, cache: RewriteCache | None = None) -> NormalForm!
\end{signature}
\begin{algorithm}[Normal-form driver]\label{alg:rewriting-nf}
\leavevmode
\begin{enumerate}
\item Require a \py!Diagram! (else \py!RewriteGrammarError!, \code!normal_form: diagram must be a Diagram, got <type>!).
\item Choose the rules: when \py!rules! is \py!None!, \py!normal_form_rules()! for a pure diagram and
  \py!mixed_normal_form_rules()! for a diagram with a discard (\cref{ch:mixed}): the forward discard
  rules, in the order of the discard axioms, followed by the normal-form rules. A given \py!rules!
  must be a non-string sequence of \py!Rule!.
\item Without a cache, create a \py!RewriteCache! whose incremental matcher tracks every keyable
  pattern of the rules (\cref{sec:rewriting-cache}).
\item Run \py!apply_until_fixpoint! on a copy of the source with the guard and cache, collecting
  every result.
\item Key the result by \py!view_key(comparison_view(outcome.diagram))!.
\end{enumerate}
\end{algorithm}

\subsection{Comparison views and keys}

\begin{definition}[Comparison view]\label{def:rewriting-view}
\py!comparison_view(diagram)! returns a new diagram in which
\begin{enumerate}
\item the ports of every Z and X spider are sorted per direction by a node-id-free key (below),
  unless the diagram carries a \bbox{}, in which case the ports stay as they are;
\item the parameter environment is empty;
\item the scalar is replaced by its simplified form, or, for a mixed diagram, by the simplified form
  of $c\cdot\bar c$, the CPM scalar~\cite{selinger2007cpm}.
\end{enumerate}
Simplification failures (\py!ScalarError!) leave the scalar unchanged. Simplified scalars are
memoized by SymPy expression, the memo cleared once it holds $4096$ entries.
\end{definition}

The sort key of a spider port $p$ with dimension text $t$ is
\begin{itemize}
\item $(0, s, i, t)$ when $p$ is a boundary entry, with $s=0$ for inputs and $1$ for outputs and $i$
  its first position in that list;
\item $(2, t, \mathit{index})$ when $p$ is unwired and not on the boundary;
\item otherwise $(1, t, g, \mathit{dir}, j, m, n, \phi, \mathit{index})$ describing the far port $q$:
  $g$ its node's generator name, $\mathit{dir}$ its direction value, $j=-1$ if its node is a Z or X
  spider and $q$'s index otherwise, $m,n$ its node's input and output counts, $\phi$ the digest of
  its node's phase or \code!"None"!, and finally $p$'s own index.
\end{itemize}
Only directions with at least two ports are sorted. The view renumbers nodes in ascending order of
the original ids and copies wires, boundary lists, scalar, and parameters through the port
permutation.

\begin{itemize}
\item \py!view_key(view)! is \py!canonical_key(view, symmetric_legs=True)! (\cref{ch:diagrams}):
  the canonical key blind to the order of a spider's ports within a direction.
\item \py!view_form(view)! is \py!canonical_form(view, symmetric_legs=True)!: the key and its
  profile.
\item \py!views_isomorphic(a, b)! is \py!isomorphic(a, b, symmetric_legs=True)!.
\item \py!same_normal_form(a, b)! is true iff both are \py!NormalForm! records (else
  \py!RewriteGrammarError!), their keys are equal, and their comparison views are isomorphic up to
  spider port order.
\item \py!algebraic_view(diagram)! is the comparison view with every phase vector and the scalar
  in the canonical algebraic form of \cref{ch:algebra} (\py!canonical_vector!,
  \py!canonical_scalar!).
\item \py!fused_isomorphic(a, b)! reduces each diagram by spider fusion and identity removal only,
  then compares comparison views up to spider port order.
\item \py!isomorphic_up_to_scalar(a, b)! compares copies with an empty parameter environment and
  scalar $1$.
\end{itemize}

The decision procedure of \cref{ch:semantics} uses these comparisons as its first rungs.

%======================================================================
\section{Equality saturation}\label{sec:rewriting-egraph}

The module \module{archytaszx.rewrite.egraph} implements a non-destructive \term{e-graph} over
whole diagrams with cost-based extraction, after equality saturation~\cite{tate2011eqsat} and the
egg system~\cite{willsey2021egg}. Unlike a term e-graph, an e-node here is an entire diagram, not
an operator applied to child classes; there is therefore no congruence closure, and two classes
merge only when a rewrite joins them. Every equality the e-graph holds is witnessed by stored
\py!apply! calls, so every extraction carries a replayable rewrite path.

\subsection{E-nodes, classes, and hash-consing}

\begin{definition}[E-node and e-class]\label{def:rewriting-enode}
An \term{e-node} is a diagram up to isomorphism of its comparison view (\cref{sec:rewriting-nf}).
E-node ids \py!ENodeId! are consecutive integers from $0$ in insertion order. Each e-node keeps the
first-seen \emph{representative} exactly as produced, its comparison view, and the profile of that
view. An \term{e-class} is a block of a union-find partition~\cite{tarjan1975unionfind} of the
e-nodes; its id \py!EClassId! is the id of its root e-node.
\end{definition}

\paragraph{Hash-consing.}
E-nodes are bucketed by view key (\py!view_form!, the canonical key up to spider port order). To find
the e-node of a view, scan the bucket of its key in insertion order: when the new view or the
stored one has a profile, the two match iff their profiles are equal; when neither has one, they
match iff \py!views_isomorphic!. A second memo maps the \emph{exact key} of a diagram (sorted nodes
with ids, sorted wire keys, both boundary lists, sorted \bboxes{}, scalar, sorted parameters) to its
e-node, so a result identical to an earlier one, node ids included, skips the view computation.

\paragraph{Union-find.}
\py!_root(i)! follows parent pointers to the root and then compresses the path, pointing every
visited node at the root. A union sets the parent of the larger root id to the smaller, so the
smaller root wins. A new child e-node is inserted with its parent pointer set to the root of its
parent's class.

\begin{signature}
\py!class EGraph!
\end{signature}
\begin{itemize}
\item \py!EGraph()!: an empty e-graph. \py!len(g)! is the e-node count.
\item \py!add(diagram) -> ENodeId!: the e-node already holding an isomorphic view, or a new root
  e-node holding a copy, appended to \py!roots!.
\item \py!lookup(diagram) -> ENodeId | None!: the e-node isomorphic under the view, or \py!None!.
\item \py!find(id)!, \py!class_of(enode)! \py!-> EClassId!: the class root.
\item \py!equivalent(a, b) -> bool!: whether the two share a class.
\item \py!members(eclass)!: every e-node of the class, ascending. \py!classes()!: every class id,
  ascending.
\item \py!diagram(enode)!: a copy of the representative.
\item \py!edges!: every stored edge in insertion order, each with a copied result diagram.
  \py!roots!: every e-node created by \py!add!, in creation order.
\end{itemize}
An id that is not an \py!int!, is a \py!bool!, or is out of range raises \py!RewriteGrammarError!
(\code!<method> must be an int id, got <type>!, \code!<method>: unknown id N!). A non-diagram
argument raises \py!RewriteGrammarError!. Every diagram, edge, and result the e-graph hands out
carries a copy of the stored diagram.

\begin{definition}[Edge]\label{def:rewriting-edge}
An \py!EGraphEdge(parent, child, rule_name, result, kind)! records one \py!apply! call on the
parent's representative; \py!result! is the \py!RewriteResult!. Its kind (\py!EdgeKind!) is
\code!TREE! when the edge created its child, whose representative is the result diagram, in the
parent's class, and \code!MERGE! when the result was view-isomorphic to an e-node in another class
and the edge unioned the two classes. A result already in the parent's class is counted as
redundant and not stored. The record validates its field types and compares by identity.
\end{definition}

\subsection{Saturation}

\begin{definition}[Saturation limits]\label{def:rewriting-limits}
\py!SaturationLimits! is a frozen record of non-negative integers (a \py!bool! or negative value
raises \py!RewriteGrammarError!):
\begin{center}\small
\begin{tabular}{@{}lrl@{}}
\toprule
Field & Default & Meaning \\
\midrule
\py!max_iterations! & $6$ & rounds per call \\
\py!max_enodes! & $512$ & e-nodes in the e-graph \\
\py!max_applications! & $20000$ & \py!apply! calls per call \\
\py!node_margin! & $4$ & added to the largest root's node count to form the node bound \\
\py!match_limit! & $1000$ & backoff threshold base \\
\py!ban_length! & $5$ & backoff ban length base \\
\bottomrule
\end{tabular}
\end{center}
\py!DEFAULT_LIMITS = SaturationLimits()!. Two module constants bound the stored size:
\py!STORED_NODE_BUDGET! $=2^{17}$ nodes across all representatives, and
\py!STORED_WEIGHT_BUDGET! $=2^{22}$, read as an e-node cap of
$\max(2, \lfloor 2^{22}/w\rfloor)$ with $w$ the largest \py!diagram_weight! among the roots, at
least $1$. \py!diagram_weight(diagram)! is the node count plus the number of explicit phase entries
over all nodes.
\end{definition}

\py!saturation_rules()! returns every rule of the registry \py!RULES! (\cref{ch:rules}) whose
\py!automatic! flag is set, in registry order.

\begin{signature}
\py!EGraph.saturate(rules=None, *, limits=DEFAULT_LIMITS, cache: RewriteCache | None = None) -> SaturationReport!
\end{signature}
\begin{algorithm}[Saturation]\label{alg:rewriting-saturate}
Let $R$ be the rules (default \py!saturation_rules()!), the node bound $\nu$ the largest node count
among the roots plus \py!node_margin!, and set the e-node cap from the weight budget.
\begin{enumerate}
\item An e-node is \emph{pending} when it was never marked expanded, or was marked under a
  different rule tuple (compared element by element by object identity), or under a smaller node
  bound than $\nu$.
\item Loop: collect the pending e-nodes in ascending id. If none, stop with \code!SATURATED!. If
  \py!max_iterations! rounds have run, stop with \code!ITERATION_LIMIT!. Start a round of the
  backoff scheduler.
\item For each pending e-node $i$ in that snapshot: the \emph{due} rules are those of $R$ not yet
  applied to $i$ in this call and allowed by the scheduler. Expand $i$ with them (below). Mark every
  due rule that the scheduler did not cut as applied to $i$; when all of $R$ is applied, mark $i$
  expanded under $(R, \nu)$.
\item End the scheduler's round.
\end{enumerate}
\emph{Expansion} of $i$, for each due rule in order and each of its matches on the representative
(through the cache's match memo when given) whose side conditions all pass:
\begin{enumerate}
\item Count a firing with the scheduler; if it refuses, abandon the rest of this rule's matches.
\item If \py!max_applications! applications were made, stop with \code!APPLICATION_LIMIT!.
\item Apply, with the representative's validation report computed once per expansion. An exception
  of one of the types \py!RewriteError!, \py!BangBoxError!, \py!DimensionError!,
  \py!GeneratorError!, \py!GraphError!, \py!PhaseError!, \py!ScalarError!, or \py!ValidateError!
  (\py!APPLICATION_ERRORS!) is recorded as \py!FailedApplication(enode, rule_name, message)! with
  message \code!<ExceptionType>: <text>!, and the loop continues.
\item Record the result: if its node count exceeds $\nu$, count it \emph{pruned}. Otherwise find
  its e-node by exact key, then by view. Not found: if the e-graph holds \py!max_enodes! e-nodes, or
  the e-node cap, or the stored nodes plus the result's would exceed \py!STORED_NODE_BUDGET!, stop
  with \code!ENODE_LIMIT!; else insert it in the parent's class with a \code!TREE! edge. Found in the
  parent's class: count it redundant. Found elsewhere: store a \code!MERGE! edge and union.
\end{enumerate}
A cache rejection raises \py!RewriteGrammarError!.
\end{algorithm}

The three limits stop the run deterministically: the e-node limit before an insertion, the
application limit before an apply, and the iteration limit before a round.

\paragraph{Backoff scheduler.}
As in egg, a rule that fires too often within one round is cut for the rest of it and banned for
a while. Per rule (by object identity) the scheduler keeps $b$, the number of bans so far, and the
round $u$ until which it is banned. A rule is allowed in round $t$ iff $u < t$. In a round, its
$c$-th firing is refused, and the rule cut, when $c > \text{\py!match_limit!}\cdot 2^{b}$; a cut
rule refuses every further firing that round. At the end of round $t$ each cut rule gets
$u = t + \text{\py!ban_length!}\cdot 2^{b}$ and $b \gets b+1$. A banned or cut rule stays unapplied
on its e-nodes, which therefore remain pending.

\begin{definition}[Saturation report]\label{def:rewriting-satreport}
\py!SaturationReport(stop_reason, iterations, enodes_added, merges, applications, redundant, pruned, failed)!
is a frozen record with field types validated. \py!SaturationStop! has the members
\code!SATURATED!, \code!ITERATION_LIMIT!, \code!ENODE_LIMIT!, and \code!APPLICATION_LIMIT!;
\py!saturated! is true iff the stop is \code!SATURATED!.
\end{definition}


\paragraph{Invariants.}
After every insertion and union the e-graph satisfies:
\begin{enumerate}
\item every e-node's representative is unique up to view isomorphism among all representatives;
\item every class id is the smallest e-node id of its class;
\item every non-root e-node has exactly one \code!TREE! edge into it, whose parent was in the same
  class when the edge was stored, so \py!path_from_root! terminates at a root created by \py!add!;
\item two e-nodes are in one class iff a chain of stored edges, read undirected, joins them;
\item every stored edge's result is the exact output of \py!apply! on its parent's representative.
\end{enumerate}

\begin{example}[A merge]\label{ex:rewriting-merge}
Add $C_3$ and then $C_2$ (\cref{ex:rewriting-chain}) to an empty e-graph: they become root e-nodes
$0$ and $1$ in classes $0$ and $1$. Saturating with \py!max_iterations=4! stops with
\code!SATURATED! after $2$ rounds, $1$ e-node added, $1$ merge, $11$ applications, and $9$
redundant results. Spider fusion on $C_3$ yields a diagram view-isomorphic to e-node $1$, stored as
the edge (0, 1, \code!spider_fusion!, \code!MERGE!), which unions the two classes into class $0$;
fusion on $C_2$ then yields the new e-node $2$ by the edge (1, 2, \code!spider_fusion!, \code!TREE!).
\py!explain(0, 1)! returns the single merge edge.
\end{example}

\subsection{Extraction and explanation}

A \py!CostFunction! maps a diagram to an \py!int!, a non-NaN \py!float!, or a flat tuple of them
(\py!bool! excluded). Two are provided: \py!node_count_cost(d)!, the node count, and
\py!default_cost(d)!, the triple (node count, wire count, number of nodes with a nonzero phase).

\begin{signature}
\py!EGraph.extract(eclass, cost=default_cost) -> Extraction!
\end{signature}
evaluates the cost on a copy of every member's representative, in ascending id, and keeps the
first strict minimum, so ties go to the lowest e-node id. A non-callable cost, a value of the wrong
type, or values that are not mutually comparable raise \py!RewriteGrammarError!. The result
\py!Extraction(eclass, enode, diagram, cost, root, results)! carries the class root, the winning
e-node, a copy of its representative, its cost, and the path from its root.

\py!EGraph.path_from_root(enode) -> (root, results)! follows \code!TREE! edges back from the e-node
to the root e-node it descends from and returns their results in forward order; replaying the
results from the root's representative reproduces the e-node's representative, node ids included.

\py!EGraph.explain(a, b) -> tuple[EGraphEdge, ...] | None! returns \py!()! when $a=b$, \py!None!
when the two are in different classes, and otherwise a shortest path of stored edges between them,
read undirected, found by breadth-first search from $a$ over the edges in insertion order. This is
an explanation in the sense of~\cite{flatt2022proofs}: each edge is one rewrite step between the
two representatives it joins (in its stored direction).

\begin{signature}
\py!simplify(diagram, rules=None, *, cost=default_cost, limits=DEFAULT_LIMITS, cache=None) -> tuple[Extraction, SaturationReport]!
\end{signature}
adds the diagram to a fresh e-graph, saturates, and extracts its class.

%======================================================================
\section{Tactics and proof search}\label{sec:rewriting-tactics}

The module \module{archytaszx.rewrite.tactics} provides a small tactic language in the style of
LCF~\cite{gordon1979lcf} and a layered, optionally bidirectional, breadth-first search that finds a
derivation between two diagrams. The derivations it finds are certified by the proof layer
(\cref{ch:proofs}); the search itself only proposes them.

\subsection{Tactics}

\begin{definition}[Tactic]\label{def:rewriting-tactic}
A \term{tactic} maps a diagram to a lazy, deterministically ordered stream of outcomes. A
\py!TacticOutcome(diagram, results)! is a frozen record, compared by identity, whose
\py!results! \emph{chain}: \py!results[0]! is \py!apply! on the input diagram, \py!results[i+1]! is
\py!apply! on \py!results[i].diagram!, and \py!diagram is results[-1].diagram!; an outcome with no
results holds the input diagram object itself. Construction raises \py!RewriteGrammarError! when
\py!diagram! is not the last result's diagram. \py!steps! returns every result's step. No tactic
mutates a diagram, so the results replay from the input id for id.
\end{definition}

The abstract class \py!Tactic! has a property \py!name!, a readable expression of the tactic (also
its \py!repr!), and the methods
\begin{itemize}
\item \py!outcomes(diagram, context=None) -> Iterator[TacticOutcome]!: every outcome, lazily, under
  the context (a fresh one when \py!None!); a non-diagram or non-context raises
  \py!RewriteGrammarError!;
\item \py!run(diagram, context=None) -> TacticOutcome | None!: the first outcome, the stream closed
  afterwards;
\item \py!all(diagram, context=None, *, limit=None) -> tuple!: the outcomes, at most \py!limit!;
\item \py!a >> b! for \py!seq(a, b)! and \py!a | b! for \py!first(a, b)!.
\end{itemize}

\begin{definition}[Tactic context]\label{def:rewriting-tcontext}
A \py!TacticContext(*, cache=None, max_applications=None)! is the mutable state shared by one run:
an optional \py!RewriteCache!, an application budget, the count \py!applications! of applications
charged, and the list \py!failures! of \py!TacticFailure(rule_name, message)! records.
\py!charge()! raises \py!BudgetExhausted(applications)!, a \py!RewriteError! with message
\code!application budget exhausted after N application(s)!, when the budget is set and reached, and
otherwise increments the count. \py!report(diagram)! returns \py!validate(diagram)! memoized per
diagram object (identity checked); \py!keep_report(diagram, report)! stores one; the memo is cleared
once it holds $4096$ entries. \py!matches(rule, diagram)! returns the rule's matches through the
cache when present (a \py!CacheError! becomes \py!RewriteGrammarError!).
\end{definition}

\paragraph{Primitives.}
\begin{itemize}
\item \py!rule(r, *, focus=None)!: \py!r! is a \py!Rule! or a registered name resolved by
  \py!lookup_rule!. For each match in order whose side conditions all pass and, with a focus, whose
  support meets the focus: charge the context, apply with the context's input report, and yield one
  outcome; an exception of \py!APPLICATION_ERRORS! (\cref{sec:rewriting-egraph}) is recorded as a
  failure and skipped. A focus must be a non-string iterable of \py!int! node ids. Name
  \code!rule(<name>)! or \code!rule(<name>, focus=[ids])!.
\item \py!identity()!: one outcome with no results. Name \code!identity!.
\item \py!fail()!: no outcome. Name \code!fail!.
\item \py!fixpoint(rules, *, guard=DEFAULT_GUARD)!: one outcome, the \py!apply_until_fixpoint! run
  over the rules (names resolved), each result charged to the context as it is produced. An
  exception of \py!APPLICATION_ERRORS! other than \py!BudgetExhausted! is recorded as a failure
  under the tactic's name, and the outcome holds the results collected so far. Name
  \code!fixpoint(r1, r2, ...)!.
\item \py!normalize(*, guard=DEFAULT_GUARD)!: \py!fixpoint! over \py!normal_form_rules()!, named
  \code!normalize!.
\end{itemize}

\paragraph{Combinators.}
\begin{itemize}
\item \py!seq(t1, ..., tn)!: depth-first bind: for every outcome of $t_1$, every outcome of
  $\mathrm{seq}(t_2,\dots,t_n)$ on its diagram, results concatenated.
\item \py!first(t1, ..., tn)!: the outcomes of the first child whose stream yields at least one.
\item \py!choice(t1, ..., tn)!: every outcome of every child, in child order.
\item \py!attempt(t)!: \py!first(t, identity())!.
\item \py!once(t)!: the first outcome of \py!t! only.
\item \py!repeat(t, *, max_times=128)!: one outcome: take the first outcome of \py!t! repeatedly, at
  most \py!max_times! times, stopping when there is none, when it has no results, or when its diagram
  is isomorphic to one already visited (bucketed by \py!canonical_key!, the step not taken).
\end{itemize}
\py!seq!, \py!first!, and \py!choice! need at least one tactic and flatten direct children of their
own kind; every argument must be a \py!Tactic!; violations raise \py!RewriteGrammarError!. The name
of a combinator is \code!label(child, ...)!, with \code!max_times=N! appended to \code!repeat! when
it is not $128$.

\paragraph{Move sets.}
\py!rewrite_tactics(rules=None)! returns one \py!rule! tactic per rule (default
\py!saturation_rules()!). \py!default_moves()! is \py!(normalize(), *rewrite_tactics())!.
\py!mixed_moves()! is a fixpoint over \py!mixed_normal_form_rules()! named
\code!normalize_mixed!, followed by one \py!rule! tactic per discard rule (\cref{ch:mixed}) and per
saturation rule.

\subsection{Proof search}

\begin{definition}[Search limits]\label{def:rewriting-slimits}
\py!SearchLimits(max_depth=4, max_states=2048, max_applications=20000, node_margin=4)! is a frozen
record of non-negative integers: moves per side, distinct states, applications, and the margin of
the node bound. \py!DEFAULT_SEARCH_LIMITS = SearchLimits()!.
\end{definition}

\py!SearchStatus! has the members \code!FOUND!, \code!EXHAUSTED!, \code!DEPTH_LIMIT!,
\code!STATE_LIMIT!, and \code!APPLICATION_LIMIT!. Failing to find a path is a status, never an
exception.

\begin{definition}[Proof path]\label{def:rewriting-ppath}
A \py!ProofStep(tactic: str, results)! is one move: the tactic's name and its nonempty chained
results. A \py!ProofPath(start, goal, forward, backward)! holds moves from \py!start! and from
\py!goal! whose final diagrams have isomorphic comparison views. Its properties are
\py!forward_results! and \py!backward_results! (concatenated), \py!forward_meet! and
\py!backward_meet! (the last result's diagram, or the start or goal), \py!length! (total result
count), \py!moves! and \py!backward_moves! (tactic names). Every diagram in a path is a copy.
\end{definition}

\py!SearchResult(status, path, states, expanded, applications, pruned, depth, failures)! has
\py!path! set exactly when \py!status! is \code!FOUND!; \py!found! tests that.

\begin{signature}
\py!search(start, goal, *, moves=None, limits=DEFAULT_SEARCH_LIMITS, bidirectional=True, cache=None) -> SearchResult!
\end{signature}
\begin{algorithm}[Layered bidirectional search]\label{alg:rewriting-search}
\leavevmode
\begin{enumerate}
\item Validate the arguments (\py!RewriteGrammarError! otherwise). The moves default to
  \py!default_moves()!. Create a context with budget \py!max_applications!. The node bound is
  $\max(|N_{\mathit{start}}|, |N_{\mathit{goal}}|) + \text{\py!node_margin!}$.
\item Insert a copy of \py!start! as a state of side \code!START! at depth $0$. If the comparison
  view of \py!goal! matches a stored state, return \code!FOUND! with an empty path. Otherwise insert a
  copy of \py!goal! as a \code!GOAL! state at depth $0$.
\item For each depth $t = 0, \dots, \text{\py!max_depth!}-1$: for each side (\code!START!, then
  \code!GOAL! when bidirectional), expand every state of that side at depth $t$ in ascending state id.
  If no state of the searched sides has depth $t+1$, stop with \code!EXHAUSTED!. After the last depth,
  stop with \code!DEPTH_LIMIT!.
\item \emph{Expanding} a state runs every move in order inside a \py!resolved_view_scope!
  (\cref{sec:rewriting-template}) and, for each outcome: skip it if it has no results; prune it if
  its node count exceeds the bound; look its comparison view up among the states. A state of the
  same side drops the outcome. A state of the other side ends the search: the forward chain is the
  chain of moves from the start root, the backward chain from the goal root, and the new move is
  appended to the side that produced it. Otherwise, if \py!max_states! states exist, stop with
  \code!STATE_LIMIT!; else store a new state one layer deeper.
\item \py!BudgetExhausted! stops with \code!APPLICATION_LIMIT!.
\end{enumerate}
\end{algorithm}
States are hash-consed as in the e-graph: bucketed by view key, and a bucket entry matches when
both profiles exist and are equal, or, when either profile is missing, when the views are
isomorphic up to spider port order. The reported \py!depth! is the depth reached: the layer index
at a meet or a limit, one more than the last expanded layer at exhaustion, and \py!max_depth! at the
depth limit. With \py!bidirectional=False! the goal state is still stored, so a forward outcome that
reaches it ends the search with an empty backward chain.

The forward results of a found path replay from \py!start! and the backward results from \py!goal!;
their final diagrams have isomorphic comparison views. The proof layer turns such a path into a
certificate by reversing the backward half (\cref{ch:proofs}).

\begin{signature}
\py!search_for(start, accept, *, moves=None, limits=DEFAULT_SEARCH_LIMITS, cache=None) -> tuple[SearchResult, AcceptedState | None]!
\end{signature}
searches forward only for the first diagram, \py!start! included, on which \py!accept! returns a
value other than \py!None!. \py!accept! is first called on the start copy; then the layers are
expanded as in \cref{alg:rewriting-search} with node bound $|N_{\mathit{start}}| + $
\py!node_margin!, and \py!accept! is called on each outcome that reaches a new state, before the
state limit is checked. A hit returns \py!AcceptedState(path, value)! with a path whose goal is a
copy of the accepted diagram and whose backward chain is empty.


\begin{example}[Chain of identities]\label{ex:rewriting-chain}
Let $C_n$ be a chain of $n$ phaseless Z spiders of dimension $3$, each with one input and one
output, wired in series, the first input and the last output on the boundary.
\begin{itemize}
\item \py!normal_form(C_3)! reaches a fixpoint after two \code!identity_removal! steps, since that
  rule precedes \code!spider_fusion! in the normal-form order; the result is a single spider, which
  identity removal leaves alone because neither of its ports carries a wire to splice through.
\item \py!simplify(C_3, limits=SaturationLimits(max_iterations=4))! saturates after $3$ rounds with
  $2$ e-nodes added, $0$ merges, $11$ applications, $9$ redundant results, and $0$ pruned; it
  extracts e-node $2$ of class $0$ at cost $(1, 0, 0)$, reached from root $0$ by two
  \code!spider_fusion! results.
\item \py!search(C_3, C_1)! finds a path at depth $0$ after expanding one state: the forward move
  \code!normalize! reaches a diagram whose view is isomorphic to the goal's, so the forward chain is
  \code!("normalize",)!, the backward chain is empty, and the path holds $2$ results.
\item \py!repeat(rule("spider_fusion"))!, named \code!repeat(rule(spider_fusion))!, takes three
  fusions on $C_4$ and ends with one node. The tactic
  \py!seq(rule("spider_fusion"), rule("identity_removal")) | normalize()! is named
  \code!first(seq(rule(spider_fusion), rule(identity_removal)), normalize)! and yields $4$ outcomes
  on $C_3$, all from its first branch.
\end{itemize}
\end{example}

%======================================================================
\section{Sheet rewrites}\label{sec:rewriting-sheet}

The module \module{archytaszx.rewrite.sheet} performs rewrites directly on a scalable diagram
(\py!ScalableDiagram!, \cref{ch:diagrams}): a diagram whose nodes sit in a tree of \emph{scales},
each a \code!COPIES! scale (its subtree is repeated $m$ times) or a \code!LEGS! scale (the ports
whose \emph{fan} it is are repeated $m$ times), with multiplicity $m$ a polynomial \py!Mult! in
multiplicity symbols, and whose boundary lists hold port references and nested \emph{bundles}, one
per \code!COPIES! scale. Sheet rewrites are the operations of the !-box calculus on such
\emph{sheets}~\cite{kissinger2014pattern}: they change the scale structure rather than a node
pattern.

Every operation is pure, refuses an input on which \py!validate_scalable! reports an issue, and
returns a frozen \py!SheetStep(operation: str, arguments: tuple[tuple[str, object], ...], before, after)!
whose arguments are normalized and sorted by name and whose \py!after! passes
\py!validate_scalable!. Before returning, each operation drops from the parameter environment the
binding of every multiplicity symbol that occurs free in \py!before! but no longer in \py!after!.
New node ids and scale ids are allocated from the largest existing id plus one, ascending.

\paragraph{Errors.}
\py!SheetError! derives from \py!RewriteError!. \py!SheetGrammarError! (also a
\py!RewriteGrammarError!) reports a malformed request: wrong type, unknown id, wrong scale kind, no
correspondence, an invalid input or result (\code!<op>: input is not valid: <issues>!,
\code!<op>: result is not valid: ...!). \py!SheetDomainError! (also a \py!RewriteDomainError!)
reports a value outside an operation's domain, such as a multiplicity that is not $1$. A scale id
must be an \py!int! naming a scale; a multiplicity argument must be a \py!Mult! or an \py!int!.

\subsection{Dissolve and kill}

\py!dissolve(s, scale)! removes a scale of multiplicity exactly $1$ (else \py!SheetDomainError!),
keeping its contents. For a \code!COPIES! scale, its nodes and child scales move to its parent,
and each of its bundles is spliced in place into the enclosing item list. For a \code!LEGS! scale,
its children move to its parent and every port whose fan it is takes the parent as fan when the
parent is a \code!LEGS! scale, and no fan otherwise.

\py!kill_scale(s, scale)! removes a scale of multiplicity exactly $0$ (else
\py!SheetDomainError!) with everything it scales. For \code!COPIES!: every node whose
\code!COPIES! chain contains the scale, their wires, the subtree's scales, and their bundles and
port references vanish; emptied bundles are removed. For \code!LEGS!: every port whose fan chain
contains the scale leaves its node, the later port indices of that direction shifting down; wires
and boundary references are renumbered accordingly; the subtree's scales vanish; emptied bundles
are removed.

\subsection{Split and join}

\begin{signature}
\py!split_scale(s, scale, first: Mult | int, *, dissolve_unit: bool = True) -> SheetStep!
\end{signature}
\begin{algorithm}[Divider]\label{alg:rewriting-split}
For a \code!COPIES! scale $c$ of multiplicity $m$ (another kind raises \py!SheetGrammarError!):
\begin{enumerate}
\item Compute $m' = m - \mathit{first}$ by expanding the difference in SymPy and reading it back as
  a \py!Mult!: every term must have a non-negative integer coefficient and factors that are symbols
  raised to positive integer powers, else \py!SheetDomainError!. If $\mathit{first} = 0$ or
  $m' = 0$, raise \py!SheetDomainError! (\code!split_scale: m = first + m' has a zero part!).
\item Create a sibling \code!COPIES! scale $c'$ of multiplicity $m'$ under $c$'s parent, with the
  next free scale id; duplicate every scale strictly below $c$ under fresh ids (in ascending id
  order), with the same kind and multiplicity and mapped parents.
\item Duplicate every node of $c$'s subtree, in ascending id order, under fresh node ids, with port
  fans and scale mapped; duplicate every wire whose first endpoint lies in the subtree.
\item Set $c$'s multiplicity to $\mathit{first}$. On each boundary side, insert directly after every
  bundle of $c$ its relabeled duplicate.
\item With \py!dissolve_unit!, dissolve each part of multiplicity exactly $1$, $c'$ first.
\end{enumerate}
\end{algorithm}
The step's arguments are \code!dissolve_unit!, \code!first! (as a \py!Mult!), and \code!scale!.

\begin{signature}
\py!join_scales(s, first, second) -> SheetStep!
\end{signature}
merges two distinct sibling \code!COPIES! scales $c_1, c_2$ into $c_1$, which carries the sum of
both multiplicities; the $c_2$ subtree is removed. The two subtrees must \emph{correspond}:
\begin{enumerate}
\item They hold equally many nodes; paired in ascending id order, each pair agrees on generator
  type, phase, and port dimensions, and on scale nesting: the scale chain from the node up to its
  root and the fan chain of every port pair have equal lengths, and their scales are paired
  position by position into a scale correspondence, which must stay a function.
\item The paired scales are exactly the two subtrees' scales; below the roots, paired scales agree
  on kind, multiplicity, and (paired) parent.
\item The wires of the first subtree, mapped through the node pairing, are exactly the wires of the
  second.
\item On each boundary side, in the direct items of the common parent, either neither bundle
  appears or the bundle of $c_2$ immediately follows that of $c_1$ and equals its relabeled image.
\end{enumerate}
Any violation raises \py!SheetGrammarError! naming it.

\subsection{Enclose}

\py!enclose(s, nodes, *, multiplicity=1)! wraps a nonempty set of node ids in a new \code!COPIES!
scale of multiplicity $1$ (another multiplicity raises \py!SheetDomainError!). Let $P$ be the scale
of the chosen nodes of least \code!COPIES! depth; all of them must share it. Every deeper chosen node
must lie in a child \code!COPIES! scale of $P$ all of whose nodes are chosen; those children are
reparented under the new scale. Every \code!LEGS! scale directly under $P$ that covers only directly
chosen nodes is reparented; one that covers some but not all raises. No wire may join a chosen node
to an unchosen one. On each side, the direct items of $P$ that are references to directly chosen
nodes or bundles of the moved children must form one contiguous run, which becomes the new bundle.

\subsection{Sheet fusion}

\begin{signature}
\py!fuse_sheet(s, wire: Wire) -> SheetStep!
\end{signature}
fuses the two same-color spiders a wire joins, once per copy of their shared scale. The wire must
belong to the diagram and not be a self-loop. Let $A$ be the endpoint node with the lower id. Both
nodes must be Z spiders or both X spiders; for X the wire must join ports of different directions;
both nodes must lie in the same scale; every surviving port, the far consumed port, and every phase
must carry exactly the dimension of $A$'s consumed port. Violations raise \py!SheetDomainError!.
The merged node gets a fresh id and the same scale; its inputs are $A$'s surviving inputs then
$B$'s, its outputs likewise, each keeping its fan; its phase is the sum of both phases (absent ones
read as zero), or none when both are absent and some port survives, and the zero vector when no
port survives. Every other wire between the pair becomes a self-loop. When the fusion rule's scalar
$s_f$ (\cref{ch:rules}) is not $1$, the diagram's scalar is multiplied by $s_f^{\,k}$, where $k$ is
the concrete size of the scale; a symbolic size raises \py!SheetDomainError!.

\subsection{Registry and replay}

\py!SHEET_OPERATIONS! maps the names \code!dissolve!, \code!enclose!, \code!fuse_sheet!,
\code!join_scales!, \code!kill_scale!, and \code!split_scale! to their functions.
\py!replay_sheet_step(step)! re-runs the named operation on \py!step.before! with the recorded
arguments and returns the fresh step; it raises \py!SheetGrammarError! for a non-step, an unknown
operation, arguments the function rejects with \py!TypeError!, recorded arguments that are not in
normalized form, or a result that differs from \py!step.after!.

%======================================================================
\section{Caches}\label{sec:rewriting-cache}

The module \module{archytaszx.rewrite.cache} memoizes pattern matches, canonical keys, and
arbitrary derived values, and re-matches a fixed pattern set incrementally after a local edit.
Every digest is the first $32$ hexadecimal characters of the SHA-256 hash of the UTF-8 \py!repr!
of a payload, so keys agree across processes. No order in this module comes from set iteration:
node and box ids are sorted, ports and wires use their sort keys, and enumeration members enter only
through their value or name. The module never imports the semantics layer.

\paragraph{Errors.}
\py!CacheError! derives from \py!Exception! (not from \py!RewriteError!).
\py!CacheGrammarError! reports an argument of the wrong type or shape; \py!CacheDomainError! a
well-formed argument naming something the diagram does not contain. The engine re-raises either as
\py!RewriteGrammarError! (\cref{sec:rewriting-engine}).

\subsection{Digests and fingerprints}

For a node $v$ of a diagram, define:
\begin{itemize}
\item its \emph{boundary slots}: the tuples (side \code!"in"!/\code!"out"!, direction value, port
  index, position in the list) of every boundary entry on $v$;
\item its \emph{membership}: for each \bbox{} in ascending id order, an entry
  \code!("node", m, c...)! if $v$ is in the node scope and \code!("port", dir, index, m, c...)! for
  each of its ports in the port scope (sorted), where $m$ is the box's multiplicity \py!repr!,
  suffixed \code!|sheets! for a sheet box, and $c\ldots$ the same strings for its ancestors,
  innermost first, the walk stopping at a repeated or missing id;
\item its \emph{incidence}: for every wire in sorted order and each endpoint on $v$, the tuple (own
  direction value, own index, far node id, far direction value, far index).
\end{itemize}
The \emph{payload} of $v$ is (generator name, input dimension \py!repr!s, output dimension
\py!repr!s, phase digest or \code!"None"!, sorted boundary slots, sorted membership).
\py!node_digest(diagram, node_id)! is the digest of the payload; \py!incidence_digest(diagram, node_id)!
the digest of the sorted incidence. Both raise \py!CacheGrammarError! for a non-diagram and
\py!CacheDomainError! for an absent node. \py!node_key(node)! is the digest of (generator name,
input and output dimension \py!repr!s, phase digest), with no ids; \py!diagram_key(diagram)! is
the fingerprint's global key.

\begin{definition}[Fingerprint]\label{def:rewriting-fingerprint}
\py!fingerprint(diagram) -> DiagramFingerprint! computes, in one pass, a frozen record
\begin{signature}
\py!DiagramFingerprint(nodes: Mapping[NodeId, str], global_key, boundary_key, env_key, boxes_key, match_key, params_key, dangling_boundary_key)!
\end{signature}
\begin{itemize}
\item \py!nodes[v]! $=\mathrm{digest}(\mathrm{digest}(\text{payload}), \mathrm{digest}(\text{sorted incidence}))$;
\item \py!boundary_key!: digest of the sort keys of both boundary lists;
\item \py!params_key!: digest of the sorted parameter items;
\item \py!env_key!: digest of (scalar \py!repr!, sorted parameter items);
\item \py!boxes_key!: digest of, per box in id order, (id, multiplicity string, ancestor chain,
  sorted node scope ids, sorted port scope keys, parent id or \py!None!);
\item \py!dangling_boundary_key!: digest of the boundary entries naming no live node, with side and
  position;
\item with body $=$ (boundary key, boxes key, ((id, entry) for every node in id order), sorted wire
  keys): \py!global_key! $=\mathrm{digest}(\text{env key}, \text{body})$ and
  \py!match_key! $=\mathrm{digest}(\text{params key}, \text{body})$, the global key with the scalar
  left out.
\end{itemize}
\end{definition}
The node entries are memoized per (node object id, boundary slots, membership, incidence), the
stored node object compared by identity; the memo is emptied once it exceeds $2^{16}$ entries.

\py!changed_nodes(other)! returns the node ids present in exactly one of the two fingerprints and
those whose entries differ. \py!global_change(other)! is true when the boundary, environment, or
box key differs. \py!match_change(other)! is true when the parameter, box, or dangling-boundary key
differs. A non-fingerprint argument raises \py!CacheGrammarError!.


\paragraph{Internal memos.}
Three further memos keep repeated scans cheap without changing any result. The incidence of every
node is computed from the diagram's incident-wire index (\py!incident_state!, \cref{ch:diagrams}):
the per-node wire ends are memoized per (node id, object identity of its sorted incident-wire
tuple), the memo emptied at $2^{17}$ entries, and the last call's result is reused for every node
the index did not touch since. Wire sort keys in the fingerprint body are read from the shared
\py!WIRE_KEYS! table. The neighbor map of the fixed-shape matchers (\cref{sec:rewriting-fixed}) is
built once per wire-set snapshot object and shared read-only; it maps only nodes that carry a wire.
The incremental matcher memoizes each held match's identity (support ids, \py!repr!) per match
object and keeps only the entries of matches it still holds after each adoption.

\begin{example}[Fingerprint and rematch]\label{ex:rewriting-fp}
For $C_2$ (\cref{ex:rewriting-chain}) the node entries are \code!7c2d2906565cda78c16fe645074803f7!
for node $0$ and \code!3f5275ba6cf65b108c6fefe4d15ef919! for node $1$; the global key is
\code!c86e64fdcaf74954b0966754a0f3cd39! and the match key \code!fcbe3cc382da44820c5b6bbe38a34ed9!.
Normalizing $C_6$ with a cache whose incremental matcher tracks the patterns of the normal-form
rules (radius $8$, the template radius) takes five \code!identity_removal! steps and six rematches,
of which only the first is a full rescan; the last rematch has $2$ dirty nodes and a closure of $2$
nodes.
\end{example}

\subsection{The LRU memo}

\py!LruMemo(max_entries, *, counters=None)! is a bounded memo over an insertion-ordered dictionary.
\py!max_entries! must be an \py!int!, not a \py!bool!, at least $1$ (else \py!CacheGrammarError!).
\begin{itemize}
\item \py!get_or_compute(key, compute)!: on a hit, move the key to the most recent end, count a hit,
  and return the value; on a miss, count a miss, compute, store at the most recent end, and evict
  least recently used entries, counting each eviction, while more than \py!max_entries! are held. A
  non-callable \py!compute! raises \py!CacheGrammarError!.
\item \py!invalidate(key)!, \py!invalidate_where(predicate)!, \py!clear()!: drop entries, counting
  each as an invalidation, and report whether or how many went.
\item \py!keys()! (least recent first), \py!len!, and \py!in! (which touches neither recency nor
  counters).
\end{itemize}
\py!CacheStats(hits, misses, evictions, invalidations, unkeyable)! is a frozen snapshot with
\py!total! $=$ hits $+$ misses and \py!hit_rate! $=$ hits$/$total, or $0.0$ with no lookups.

\subsection{Pattern keys and the match cache}

\py!pattern_key(pattern)! is \code!<module>.<qualname>|<repr>!. When the \py!repr! contains
\code! object at 0x! (an identity-dependent \py!repr!), the key is \code!<module>.<qualname>|! if
the pattern carries no state, and the call raises \py!CacheGrammarError! otherwise. A pattern carries
state when it has an instance dictionary entry, a filled slot in any class of its method resolution
order, or a class attribute (outside \py!object! and \py!Pattern!) bound to a list, dictionary, set,
or byte array. Frozen dataclass patterns therefore key by value; \py!TemplatePattern! keys by its
custom \py!repr!. \py!keyable(pattern)! is whether \py!pattern_key! succeeds.

\py!MatchCache(max_entries=256)! holds two LRU memos sharing one set of counters:
\begin{itemize}
\item \py!matches(pattern, diagram, *, fingerprint=None)!: the matches memoized on
  (pattern key, match key); an unkeyable pattern is scanned without caching and counted as
  \py!unkeyable!. Since the match key omits the scalar, diagrams that differ only in the scalar share
  matches.
\item \py!matches_anchored(pattern, diagram, anchors, *, fingerprint=None)!: the anchored scan,
  never cached; anchors must be a set of non-\py!bool! integers.
\item \py!canonical(diagram, *, fingerprint=None)!: \py!canonical_key(diagram)! memoized on the
  global key.
\item \py!invalidate(fingerprint)!: drop the matches under its match key and the canonical key under
  its global key. \py!clear()!: drop everything.
\end{itemize}
A missing fingerprint is computed on the spot.

\py!RewriteCache(match_cache=MatchCache(), incremental: IncrementalMatcher | None = None)! is the
one cache object the engine accepts; it delegates \py!matches!, \py!canonical!, \py!fingerprint!,
and \py!stats! to the match cache, and \py!clear()! also resets the incremental baseline.

\subsection{Value memo}

\py!ValueMemo(compute, key, max_entries=1024)! memoizes \py!compute(subject)! under the string
\py!key(subject)! (a non-string key raises \py!CacheGrammarError!). A \tool{NumPy}{https://numpy.org}
array~\cite{harris2020numpy} is stored as a private copy and handed out as a read-only view; other
values pass unchanged. It offers \py!get!, \py!key_for!, \py!invalidate!, \py!clear!, \py!stats!,
\py!max_entries!, and \py!len!. The semantics layer (\cref{ch:semantics}) instantiates it for
denotations.

\subsection{Incremental rematching}

\py!IncrementalMatcher(patterns, *, cache: MatchCache | None = None)! re-matches a fixed sequence of
keyable patterns after a local edit, rescanning only the affected region. Its \py!radius! is the
largest \py!locality_radius! among the patterns ($0$ for none). It keeps a \emph{baseline}: a
fingerprint, the node adjacency, the box scopes (per box, the descriptor of \py!boxes_key! without
the id), and the matches per pattern key.
\py!seed(diagram)! scans every pattern in full (through the match cache when given) and adopts the
result. \py!reset()! drops the baseline; \py!clear()! also zeroes the counters.

\begin{algorithm}[Incremental rematch]\label{alg:rewriting-rematch}
\py!rematch(diagram, *, fingerprint=None)!:
\begin{enumerate}
\item Without a baseline, or when \py!match_change! holds against it, scan every pattern in full,
  adopt the result, and count a full rescan.
\item The \emph{dirty} set is the changed nodes of the fingerprint, plus every node scoped (by node
  or port scope) by a box whose descriptor is new or changed, plus the node scope of every old box
  whose descriptor disappeared or changed.
\item For every radius $r$ in use, the closure $C_r$ is the dirty set grown by $r$ breadth-first hops
  over the union of the old and the new adjacency.
\item For each pattern with radius $r$: a baseline match is \emph{retained} when its support is
  disjoint from $C_r$ and every support node has the same fingerprint entry as in the baseline. The
  fresh matches are \py!find_matches_anchored(diagram, C_r)!. Retained and fresh matches are merged,
  deduplicated by identity (support ids, \py!repr!), retained first, and sorted by
  (\py!order_key!, \py!repr!). When some baseline or fresh match has no integer support or an
  identity-dependent \py!repr!, or the pattern raises any exception other than \py!CacheError!, that
  pattern is scanned in full instead.
\item Adopt the result as the new baseline and update the counters.
\end{enumerate}
\end{algorithm}
\py!IncrementalStats(dirty_nodes, closure_nodes, retained, rescanned, full_rescans, rematches)!
records the sizes of the last rematch and the running totals of full rescans and rematches.

The normal-form driver equips every run with an incremental matcher over its rules' keyable
patterns (\cref{sec:rewriting-nf}); the fixpoint strategy then rematches once per step, so a local
rewrite rescans only the neighborhood of the nodes it changed.

%======================================================================
\section{Quick reference}\label{sec:rewriting-quickref}

\begin{xltabular}{\textwidth}{@{}>{\raggedright\arraybackslash}p{0.45\textwidth}>{\raggedright\arraybackslash}X@{}}
\caption{Public names of the rewriting layer and where they are specified.}\label{tab:rewriting-quickref}\\
\toprule
Name & Specified in \\
\midrule
\endfirsthead
\toprule
Name & Specified in \\
\midrule
\endhead
\multicolumn{2}{@{}l}{\file{rule.py}}\\
\quad\py!RewriteError!, \py!RewriteDomainError!, \py!RewriteGrammarError! & \Cref{ch:rewriting}, Errors \\
\quad\py!SideCondition!, \py!SideConditionOutcome!, \py!check_side_condition_coverage! & \Cref{sec:rewriting-rule} \\
\quad\py!ConstraintSource!, \py!ConstraintOutcome!, \py!DimensionConstraint!, \py!assumed! & \Cref{def:rewriting-csource,def:rewriting-dconstraint} \\
\quad\py!DimensionGuard!, \py!DimensionGuardKind!, \py!GuardOutcome! & \Cref{def:rewriting-guard} \\
\quad\py!Quantifiers!, \py!Match!, \py!Pattern!, \py!BuildResult!, \py!RuleBuilder!, \py!Rule! & \Cref{def:rewriting-match,def:rewriting-pattern,def:rewriting-buildresult,def:rewriting-rulerecord} \\
\multicolumn{2}{@{}l}{\file{match.py}}\\
\quad\py!FUSION_SIDE_CONDITIONS!, \py!FusionMatch!, \py!FusionResolution!, \py!resolve_fusion_match!, \py!reattach_phase!, \py!find_matches!, \py!FusionPattern!, \py!innermost_node_scope_box! & \Cref{sec:rewriting-fusion} \\
\quad\py!FOURIER_CHAIN_LENGTH!, \py!FourierMatch!, \py!find_fourier_matches!, \py!FourierCancellationPattern! & \Cref{sec:rewriting-fixed} \\
\quad\py!colour_pair!, \py!CapMatch!, \py!find_cap_matches!, \py!CapPattern! & \Cref{sec:rewriting-fixed} \\
\quad\py!IdentityMatch!, \py!find_identity_matches!, \py!IdentityRemovalPattern! & \Cref{sec:rewriting-identity} \\
\quad\py!TriangleMatch!, \py!StateCopyMatch!, \py!HopfMatch!, \py!BialgebraMatch!, \py!FourierStateMatch! and their finders and patterns & \Cref{sec:rewriting-fixed} \\
\multicolumn{2}{@{}l}{\file{match_template.py}}\\
\quad\py!Param!, \py!Fixed!, \py!PHASELESS!, \py!Var!, \py!Derived!, \py!Kj!, \py!k_vector!, \py!kj_index! & \Cref{sec:rewriting-template} \\
\quad\py!TNode!, \py!Template!, \py!Side!, \py!Condition!, \py!TemplateAxiom!, \py!TemplateTooLarge!, \py!check_template_budget! & \Cref{def:rewriting-template,def:rewriting-axiom} \\
\quad\py!match_template!, \py!TemplatePattern!, \py!TemplateMatch!, \py!template_scalar!, \py!template_builder! & \Cref{alg:rewriting-subgraph,alg:rewriting-templatefind,alg:rewriting-rebuild} \\
\quad\py!CONCRETE_AXIOMS!, \py!resolved_view!, \py!resolved_view_scope!, \py!legs_of!, \py!neighbours!, \py!ranges! & \Cref{sec:rewriting-template} \\
\multicolumn{2}{@{}l}{\file{match_zxw.py}, \file{match_box.py}}\\
\quad W, triangle, connective, and port-box matchers & \Cref{sec:rewriting-zxw,sec:rewriting-box} \\
\multicolumn{2}{@{}l}{\file{engine.py}}\\
\quad\py!apply!, \py!RewriteStep!, \py!RewriteResult! & \Cref{alg:rewriting-apply,def:rewriting-step} \\
\quad\py!StopReason!, \py!TerminationGuard!, \py!DEFAULT_GUARD!, \py!TerminationGuardTripped!, \py!StrategyOutcome!, \py!apply_until_fixpoint! & \Cref{alg:rewriting-fixpoint} \\
\quad\py!NORMAL_FORM_RULE_NAMES!, \py!normal_form_rules!, \py!missing_normal_form_rule_names!, \py!toward_normal_form! & \Cref{sec:rewriting-engine} \\
\multicolumn{2}{@{}l}{\file{normal_form.py}}\\
\quad\py!NormalForm!, \py!normal_form!, \py!mixed_normal_form_rules!, \py!comparison_view!, \py!algebraic_view!, \py!view_key!, \py!view_form!, \py!views_isomorphic!, \py!same_normal_form!, \py!fused_isomorphic!, \py!isomorphic_up_to_scalar! & \Cref{sec:rewriting-nf} \\
\multicolumn{2}{@{}l}{\file{egraph.py}}\\
\quad\py!EGraph!, \py!EGraphEdge!, \py!EdgeKind!, \py!SaturationLimits!, \py!DEFAULT_LIMITS!, \py!SaturationStop!, \py!SaturationReport!, \py!FailedApplication!, \py!Extraction!, \py!saturation_rules!, \py!diagram_weight!, \py!node_count_cost!, \py!default_cost!, \py!simplify!, \py!APPLICATION_ERRORS!, \py!STORED_NODE_BUDGET!, \py!STORED_WEIGHT_BUDGET! & \Cref{sec:rewriting-egraph} \\
\multicolumn{2}{@{}l}{\file{tactics.py}}\\
\quad\py!Tactic!, \py!TacticOutcome!, \py!TacticContext!, \py!TacticFailure!, \py!BudgetExhausted!, \py!rule!, \py!identity!, \py!fail!, \py!seq!, \py!first!, \py!choice!, \py!attempt!, \py!once!, \py!repeat!, \py!fixpoint!, \py!normalize!, \py!rewrite_tactics!, \py!default_moves!, \py!mixed_moves! & \Cref{sec:rewriting-tactics} \\
\quad\py!search!, \py!search_for!, \py!SearchLimits!, \py!SearchStatus!, \py!SearchResult!, \py!ProofStep!, \py!ProofPath!, \py!AcceptedState!, \py!Side! & \Cref{alg:rewriting-search} \\
\multicolumn{2}{@{}l}{\file{sheet.py}}\\
\quad\py!SheetError!, \py!SheetGrammarError!, \py!SheetDomainError!, \py!SheetStep!, \py!dissolve!, \py!kill_scale!, \py!split_scale!, \py!join_scales!, \py!enclose!, \py!fuse_sheet!, \py!SHEET_OPERATIONS!, \py!replay_sheet_step! & \Cref{sec:rewriting-sheet} \\
\multicolumn{2}{@{}l}{\file{cache.py}}\\
\quad\py!CacheError!, \py!node_digest!, \py!incidence_digest!, \py!fingerprint!, \py!DiagramFingerprint!, \py!node_key!, \py!diagram_key!, \py!LruMemo!, \py!CacheStats!, \py!pattern_key!, \py!keyable!, \py!MatchCache!, \py!ValueMemo!, \py!IncrementalMatcher!, \py!IncrementalStats!, \py!RewriteCache! & \Cref{sec:rewriting-cache} \\
\bottomrule
\end{xltabular}
